\begin{center}
{\Large\bfseries One normalization for the original theta comparisons}\\[4pt]
Edition 0.74 \quad 24 November 2026
\end{center}

The theta lift has two readings. One integrates over a definite
unitary group. The other evaluates a section on a marked moduli
space. To transfer rationality from the base tuple to every tuple,
the two readings must agree up to a factor chosen before the tuple,
its coefficient field and the marking vary. We construct that
factor, track the lattice index in the comparison, and determine
which cusp sections the resulting argument already generates.
The final chapter explains how the base identities descend together
to the bundle selected by a linked family.

\begingroup\emergencystretch=2em
\section{The marked higher dictionary on the prescribed coefficient field}
\label{sec:prescribed-higher-dictionary}

The higher marking changes the integral period lattices. To read the
character on the resulting moduli scheme one needs a period chart
for each adelic lattice, using the given embedding of the lower
lattice and the given archimedean frame. We first supply the two
section laws used to normalize those charts.

% BEGIN D74U1
\paragraph{D74U1. Sections of a polarized geometric fiber.}
\label{claim:D74U1}
Let $P$ be an original $\mathrm{PolAV}$ over a scheme $S$. On every
geometric fiber $B$ of dimension $g=P.g$, write $L$ for the pulled
line. For every $q\geq1$ and $c\in B(k)$, both $L^q$ and
$N=L^{q-1}\otimes t_c^*L$ have $q^g$ independent global sections
forming a basis over the algebraically closed residue field $k$, and
their higher cohomology vanishes.

\begin{proof}
PBC25 transports the polarization to the fiber, including the field
$\mathrm{KTriv}$. In particular $L$ is ample and symmetric.
Consider the polarization isogeny
$\lambda_L:B\to B^\vee$, $a\mapsto[t_a^*L\otimes L^{-1}]$.
We check its scheme-theoretic kernel, which matters in positive
characteristic. For every affine $T/k$ and $a\in\ker\lambda_L(T)$,
put $D_a=t_a^*L_T\otimes L_T^{-1}$. Rigidify $D_a$ by tensoring
with the inverse of the pullback of its identity restriction.
The rigidified bundle then has zero relative Picard class. It has a
unique rigidified trivialization fppf locally on $T$: membership
in the zero class gives a pullback from the base locally, and
rigidifying such a pullback trivializes it. Uniqueness follows
from $(B_T\to T)_*\mathcal O_{B_T}=\mathcal O_T$. For this last
identity over arbitrary $T$, one may work over an affine of $T$
and use flat base change from the field $k$ and
$H^0(B,\mathcal O_B)=k$. Uniqueness descends the trivialization.
Thus $D_a$ is the pullback of a line on $T$ and has the generators
Zariski locally on $T$ required by $\mathrm{KTriv}$. The latter
gives $a=0$. Affine tests include the possibly nonreduced kernel
itself. The isogeny has degree one.

The degree formula and Riemann--Roch give
$\deg\lambda_L=\chi(L)^2$ and $\chi(L^q)=q^g\chi(L)$.
Ampleness gives positive top self-intersection in positive dimension,
so $\chi(L)=1$. In dimension zero $B=\operatorname{Spec} k$
and the formula holds directly. For vanishing, choose arbitrarily
large integers $b$ invertible in $k$. The pullback
$[b]^*L^q\simeq L^{qb^2}$ follows from symmetry and the theorem of
the cube. The unit followed by the trace for $[b]$ multiplies by
the invertible degree $b^{2g}$. After applying the projection formula,
$H^j(B,L^q)$ is a summand of $H^j(B,L^{qb^2})$. Serre vanishing
therefore gives the assertion for $j>0$.

Choose $a\in B(k)$ with $qa=c$, using surjectivity of multiplication
on the abelian variety, also when the characteristic divides $q$.
The theorem of the square gives
$\lambda_{L^q}(a)=q\lambda_L(a)=\lambda_L(c)$, hence
$N\simeq t_a^*L^q$. Translation preserves cohomology. This choice
is needed only on an algebraically closed fiber.
The classical degree and Riemann--Roch formulas are
J.~S.~Milne, \emph{Abelian Varieties}, Theorem~13.3(a,b),
\url{https://www.jmilne.org/math/xnotes/AVs.pdf}.
\end{proof}
% END D74U1

% BEGIN D74U2
\paragraph{D74U2. The section bases before every affine test.}
\label{claim:D74U2}
For the actual polarized data the geometric rank,
universal section bases and residue lifts satisfy all four native laws:
\[
 \begin{gathered}
 \mathrm{GlobalSectionsRankAlgClosed},\qquad \mathrm{SecBaseChangeAnySt},\\
 \mathrm{SecBaseChangeSt},\qquad \mathrm{SecTwistSt}.
 \end{gathered}
\]

\begin{proof}
The absolute rank is D74U1 with the identity pullback and scalar
identifications of SGV2. For the bases let $p:B\to S'$ be the
original abelian scheme, or its pullback to the affine local test
of the twist law. Write $N$ for the indicated power or translated
power there. The map $p$ is proper, flat and of finite presentation
by the abelian-scheme fields. The line $N$ is of finite presentation
and flat over $S'$. The proper-support perfect cohomology theorem
makes $Rp_*N$ perfect with formation commuting with arbitrary derived
base change (Stacks Project, Lemma 36.30.7,
\url{https://stacks.math.columbia.edu/tag/0A1J}).
In particular its derived geometric fibers are the complexes computed
in D74U1.

At any point choose a finite free complex representing $Rp_*N$
locally. Split each unit pivot of each differential at the point as
a two-term identity summand, shrinking to invert its pivot. In the
remaining complex all differentials are zero modulo the maximal ideal.
After passage to the algebraic closure of the residue field, D74U1
therefore identifies the rank in each degree with the dimension of
the cohomology in that degree. Only degree zero remains and its rank
is $q^g$. The construction takes finitely many pivots, so gives one
open neighborhood and a basis there before choosing any later
affine test. In degree zero the derived comparison is the usual
pullback of sections, by functoriality. Thus on any affine test the
pulled basis spans and is independent, including over a nonflat
algebra. By the chart gluing in GSL1, every list satisfying the
original $\mathrm{IsPullSec}$ consists of these same pullbacks.
This proves the arbitrary-scheme law and its locally Noetherian
restriction, for every positive power.

For the translated line on $\operatorname{Spec}R$ with $R$ local,
the resulting finite projective module is free of rank $q^g$. Its
basis gives unique coordinates for each section on each of the
source's common charts, precisely $\mathrm{IsSecBasis}$. An arbitrary
section on the residue field has a coordinate vector there by base
change. Lift each of its $q^g$ coordinates to $R$ and form that
linear combination over $R$. Its pullback has the required
coordinates on every common chart by GSL1, giving the specified
$\mathrm{IsPullSecLC}$ equality. This proves both clauses of the
twist law.
\end{proof}
% END D74U2

% BEGIN D74U3
\paragraph{D74U3. Upper period models with the specified lower frame.}
\label{claim:D74U3}
Fix a slot $i$, put $V=W_c(D,i)$, $r=D.R.r$ and $D_w=T.Q.Dw(i+1)$.
Take the prescribed level, upper moduli object and coefficient record
of A73H2 and L73D2. Write $M=c'.F=c_2.F$ as an intermediate field
of $\mathbb C$. Use the original $\iota=\mathrm{iotaW}(T,i)$ and
$x_0=T.x_0(i+1,\mathrm{sOv}(T,c',\mathrm{hle}))$. There is a finite
period family $\mathcal F$ and an actual upper $U_2$ satisfying
$\mathrm{IsUnf}$, with period indices $\mathrm{ik}(k)=k$ and
integral identifications
$\phi_k:\mathcal F.D_k.\mathrm{La}\simeq_{\mathcal O_E}\Lambda_k$.
Every chart satisfies the native $\mathrm{ModelW}$ and the
period, exponential and derivative clauses of $\mathrm{DictLat}$.

\begin{proof}
The coefficient equivalence $e_F$ between the old and upper field
is induced by their equality inside $\mathbb C$. Thus the embedding
of $c_2.F$ given by $(\mathrm{inC}\ c')\circ e_F^{-1}$ is precisely
the inclusion. Its restriction to the level field has the required
root value. Choose a bijection $e$ from $\{1,\ldots,g\}$ to the
pairs $(\tau,a)$ with $a$ ranging over $D_w.\mathrm{mult}(\tau)$
indices. Use the source relabeling for this good embedding and $e$. This gives
the required $\mathrm{IsTypeFn}$ since each relabeled embedding
occurs with its prescribed multiplicity.
The adelic pins, space, uniformization, model and orthogonality laws in $\mathrm{LawsJ}$ give positive
$p+q$, rank $r$, the integral isomorphism $\iota$ onto the identity
lattice in $V$, its alternating trace form, and the given orthogonal
frame $x_0$ at this embedding. These are the fixed
$\mathrm{IotaOK}$ and frame fields used at entry to the cover
construction.

HUF21 applied to this upper moduli scheme gives $\mathcal F$.
Use PUN11 with $G=-V.\mathrm{gram}^{\mathsf T}$,
$\xi_0=-D.\mathrm{host}.\xi_0$, and $\iota$.
In these coordinates the source alternating form is
$\operatorname{Tr}_{E/\mathbb Q}(-\xi_0\mathrm{herm}_G(\iota a,\iota b))$:
interchange the arguments in the lower trace formula and use
alternation, Hermitian symmetry and invariance of trace under
conjugation. The image spans since it is the full identity lattice.
The chosen $\xi_0$ has negative imaginary part on the host CM type.
PUN11 yields on every component the displayed $\phi_k$, the literal
$G,\xi_0$, and positive-sign index $p_0=D_w.q$. Hence the
signature already prescribed for $x_0$ applies on every chart.

For each CM embedding $\tau$ use the homomorphism
$\theta_\tau:\mathbb A_{F,\infty}\otimes_F E\to\mathbb C$
given on elementary tensors by projection to the real place
$\tau|_F$ followed by multiplication by $\tau$ on the second
factor. Its restriction to the first factor is continuous and to the second is $\tau$, giving the original $\mathrm{IsInfEmb}$ restrictions.
Use its restriction to the CM type of each chart as $\mathrm{th}_k$.
Apply NFL17 with D74U2 to obtain a single $\mathrm{FramesOK}$ family
on these models. NFL21 gives $\mathrm{IsUnf}$ for its $U_2$. To
identify the first half of $\mathrm{DictLat}$, the trace identity on
$\phi_k$ is the $\mathrm{ModelBase}$ pairing field. For every
$x\in (U_2.X_k).\mathrm{dom}$ take the period homomorphism with
coordinates $\mathrm{perFr}(U_2.X_k,x,\phi_k a,\mathrm{relab}(e_j)_1,
(e_j)_2)$. The first field of $\mathrm{FramesOK}$ gives its
$\mathrm{IsPeriod}$, the equality of the point with $\mathcal F.u_k(P)$,
and the derivative of $\mathcal F.\mathrm{ex}_k(P)$ on each individual
unit vector indexed by $e^{-1}(\tau,a)$. These are the original
chart identities for all frames in the indicated domain.
\end{proof}
% END D74U3

% BEGIN D74U4
\paragraph{D74U4. A correctly marked chart for every adelic element.}
\label{claim:D74U4}
With the models of D74U3 there is a single dictionary
$h\mapsto(k(h),x(h))$ on $G(\mathbb A_F)$ satisfying the remaining
clauses of $\mathrm{DictLat}$, and every $x(h)$ lies in its indicated
chart domain.

\begin{proof}
Write $O=\mathcal O_E$, $\Lambda=\iota D_w.\mathrm{La}$, and
$N=l_2.M$. For $h\in G(\mathbb A_F)$ take
\[
 \Lambda_h= E^r\cap h_f\widehat\Lambda.
\]
This is a full $O$-lattice with completion $h_f\widehat\Lambda$.
Indeed the two completed lattices are commensurable, agreeing away
from finitely many primes. Choose a positive integer $a$ with
$a\widehat\Lambda\subset h_f\widehat\Lambda\subset
a^{-1}\widehat\Lambda$. The Chinese remainder isomorphism for
$a^{-1}\Lambda/a\Lambda$ identifies its global quotient with the
product of its local quotients. Taking the inverse image of
$h_f\widehat\Lambda/a\widehat\Lambda$ proves the assertion.
The trace pairing of D74U3 is perfect on $\Lambda_h$: at each
rational prime this follows by transport of the perfect pairing
on $\Lambda$ by the unitary $h_f$, and the rational integral dual
is the intersection of its local duals. The alternating pairing,
$O$-action, dimension and signature now define a datum $D_h$ on
$\Lambda_h$. In detail the group is finite free over $\mathbb Z$,
the relation between multiplication by $b$ and by $\bar b$ follows
from the Hermitian pairing, and its real span is the same as that
of $\Lambda$. Hence the original compatible complex structure on
this real vector space has the same positivity and Lie-type traces
on $D_h$.

At every $s>0$ the map $h_f$ induces
$\Lambda/s\Lambda\simeq\Lambda_h/s\Lambda_h$ by completion and
reduction. Choose an $O$-linear lift
$f_s:\Lambda\to\Lambda_h$; it exists since a finite torsion-free
module over the Dedekind domain $O$ is projective. At the unused
index zero set $f_0=0$. For all positive $s,u$ the reductions of
$f_{su}$ and $f_s$ modulo $s$ agree since both reduce $h_f$.
Their pairings agree with the original alternating pairing modulo
$s$ since $h_f$ is unitary, so all pairing multipliers can be one.
The root requirement at auxiliary level $A_2=1$ holds. The combined
modulus is $l_2.M'=\operatorname{lcm}(N,1)=N$, hence the specified
marking is $t_h=f_N\iota$. These data put $(D_h,t_h)$ in the original marked
genus.

By HUF14 and the choice of the family in HUF21, this class has
a representative $\mathcal F.D_k$ and an integral $O$-linear
isometry $u:\Lambda_h\to\mathcal F.D_k.\mathrm{La}$ with
\[
 u f_N(\iota a)-\mathcal F.t_k(a)\ \in N\mathcal F.D_k.\mathrm{La}
 \qquad(a\in\overline{\mathfrak b}_2D_w.\mathrm{La}).
\]
The second marking congruence modulo $A_2=1$ is automatic.
Equivalently one may use PUN15 and PUN22 to construct the point for
this marked class and PUN23 to obtain $u$ from the equality with
its family point. The map $\phi_k u$ extends uniquely to an
$E$-linear automorphism $\gamma$ of $E^r$. It preserves the
Hermitian form: apply equality of trace pairings to $ba,b'$ for
every $b\in O$, then use the perfect field trace pairing and
$\xi_0\ne0$. This puts $\gamma$ in the rational unitary group.
It sends $\Lambda_h$ onto $\Lambda_k$, hence
$\Lambda_k=\Lambda_{\gamma h}$, using rationality of $\gamma$.
The displayed congruence and
$f_N(\iota a)\equiv h_f\iota a\pmod{N\widehat\Lambda_h}$ give
the exact finite part
\[
 \phi_k(\mathcal F.t_k(a))-(\gamma h)_f\iota a\ \in
             N\widehat\Lambda_{\gamma h}
 \qquad(a\in\overline{\mathfrak b}_2D_w.\mathrm{La}).
\]
All rational vectors in this congruence use their finite diagonal
embeddings.

Choose such $k,\gamma$ for each $h$, and set
$x(h)_\tau=\operatorname{archM}_{\theta_\tau}(\gamma h)x_0(\tau)$.
The first factor preserves the chart Gram matrix, so the Gram matrix
of $x(h)_\tau$ is the same signed diagonal as that of $x_0(\tau)$.
Its leading principal minors consequently satisfy all the signed
positivity conditions of the chart, and its determinant is nonzero.
This proves domain membership. The chosen $\gamma$ now supplies
the rational witness, the lattice equality, the finite congruence
and the archimedean equality for this $h$, with the rank
identification $V.\mathrm{rank}=r$ throughout. Together with D74U3
these are all the fields of the precise $\mathrm{DictLat}$.
\end{proof}
% END D74U4

% BEGIN D74U5
\paragraph{D74U5. The original all-tuple cover packet over the prescribed field.}
\label{claim:D74U5}
Fix an original $T,R$, their $\mathrm{LawsJ}$ and level-character
laws, $R.\mathrm{fam\_roots}$, $\mathrm{XiVal}(R)$ and
$\mathrm{CoverField}(T,R)$. For every $i,n$, take any prescribed
$\{I_t,l_{2,t}\}_t$ with the original $\mathrm{IsCoverLevel}$ for
$d_t=\mathrm{tupPrime}(R,t)$ and $\mathrm{HalfIs}$ for the given lower
half ideal. For every coefficient record $c'$ with
$T.Q.c.F\subseteq c'.F=M=\mathrm{embLL}(L)$, these data give
\[
       \mathrm{CoversAt}(T,R,i,n,\mathrm{tupPrime}(R),I,l_2,c',
                              \mathrm{hle}).
\]
On the RTG1 and RTG10 inputs the resulting packet gives the original
$\mathrm{PinsOverL}$ by RTG12.

\begin{proof}
At each $t$, the level fields give $N=md_t$, $(d_t,m)=1$,
$\mathfrak b_2=\mathfrak b I_t$, $I_t+\bar I_t=O$ and $A_2=1$.
The lower auxiliary level divides $A_2$ and is one. From
$\mathfrak b_2\bar{\mathfrak b}_2=(N)$ and
$\mathfrak b\bar{\mathfrak b}=(m)$, cancellation of the nonzero
principal ideal $(m)$ gives $I_t\bar I_t=(d_t)$.
Thus the relative construction A73H1--A73H2 applies with exactly
these same levels. In L73D2 its upper coefficient field can be $M$:
the old coefficient record contains a primitive $m$th root and
$R.\mathrm{fam\_roots}$ gives one for each distinct prime in
$d_t$. Products at coprime orders give a primitive $N$th root.
All its powers belong to $M$, including the specified complex root.
The auxiliary level is one and the original quadratic base is already
in $c'$. The field identification on the actual upper and lower
coefficient records is the identity inside $\mathbb C$. L73D2 now
gives the original $\Delta_t$ at exponent $n$, its level map, framed
lifting, initial character rule and canonical connection.

For each marked prime $P\in R.Q(t)$, the $\mathrm{HalfIs}$ condition
gives $I_t\subseteq P$. The relation $(N)=\mathfrak b_2
\bar{\mathfrak b}_2$ implies $N\in\mathfrak b_2\subseteq I_t$.
Thus the same ramified-place guard as in T73K1 holds for these
prescribed levels as well. Choose its character $\xi_t$ and positive
order bound $m_t$ for the determinants. For any adelic $h$ its
determinant root exists by local Hilbert 90 as in T73K1, so its
value lies in $M$ by $\mathrm{XiVal}$. The arithmetic hypothesis
$\mathrm{CoverField}$ supplies $\mathrm{FnAlgClosed}$ for the
particular constructed upper scheme on this field. On the models and
dictionary of D74U3--D74U4 apply T73K2--T73K3. They produce a
higher rule $f_t$ with $f_t^{m_t}=1$, preserving the full
$\mathrm{IsLev}$ on replacing $\Delta_t.c_\delta$ by $f_t$. In
particular the value formula
\[
 \mathrm{rdU}(U_{2,t}, f_t,\mathrm{dict}_{2,t}(h))
      = \xi_t(\det((\gamma h)_f))
\]
holds for \emph{every} $\gamma$ satisfying
$\mathrm{MarkBy}(h,\gamma)$, including those other than the witness
chosen to define the frame.

Use these tuplewise choices as functions of $t$. The prescribed
integer, ideal and upper level are exactly those of $\Delta_t$.
The level law, $A_2=1$ and the star-ring connection witness give
$\mathrm{CovPin}$. The original half condition and T73K1 give
$\mathrm{HalfIs}$ and $\mathrm{XiIs}$. Take $c_0=1$, of order one,
before testing any tuple in $\mathrm{CXiVal}$. The uniformizations,
component models, index maps, integral isomorphisms, infinite
embeddings and marked formulas just proved give all its fields.
The repeated model clause follows from D74U3, and the last
positive-order clause uses $m_t$. This construction works also when
$d_t=1$ or $n=0$. On the rational-test input RTG12 uses these
same levels, coefficient record and lower frame to fill its eleven
pin fields. The component-field and character-value assumptions in
the displayed cover implication are statements about the original
outer arithmetic data.
\end{proof}
% END D74U5

\section{A common normalization of the point readings}

Fix a rank index $i$ and write $V=V_c(D,i)$, $W=W_c(D,i)$,
$n=\operatorname{rank} V$ and $R$ for the fixed real datum. Use the
original core, its weight and vector ties, and the extension of its
joint-harmonic map (\texttt{IoExt}). The tuple set, the coefficient
field of a tuple record $O$, and its class representatives may vary.
The coefficient map on the original lower field is fixed. The issue
is to choose a function of the automorphic point before any of these
records have been specified.

% BEGIN P74I1
\paragraph{P74I1. Comparing two projectively equivariant readings.}
Let $M$ be a nonzero irreducible finite-dimensional complex
representation of a group $\Gamma$, generated by elements with
$\operatorname{tr}\rho(\gamma)\ne0$. Let $A,B:M\to N$ be linear.
Suppose $A\rho(\gamma)=S_\gamma A$ and
$B\rho(\gamma)=c_\gamma S_\gamma B$, with linear $S_\gamma$
on $N$. Suppose also that $Aw=Bz\ne0$. Both kernels are invariant
and therefore zero. Since $\rho(\gamma)$ is invertible,
$c_\gamma\ne0$. The subspace $A^{-1}(B(M))$ contains $w$ and
is invariant since $S_\gamma B(y)=B(c_\gamma^{-1}\rho(\gamma)y)$.
It follows that $T=B^{-1}A$ is defined on all of $M$ and is invertible.
Conjugation by $T$ carries $\rho(\gamma)$ to
$c_\gamma^{-1}\rho(\gamma)$. Taking traces at the stated
generators gives $c_\gamma=1$ there. Commutation with $T$ extends
to their products and inverses. Any nonzero eigenspace of $T$ is now
invariant, so $T=aI$. In particular $A=aB$ and $aw=z$.

Here $\Gamma=U(V)(F)$ acts through the compact archimedean group of
the definite space $V$. Its dense image can be seen in the Cayley
chart at the identity. The skew-Hermitian Lie algebra is an $F$-vector
space and weak approximation gives simultaneous rational
approximation at the archimedean places. The chart
$X\mapsto(1-X)(1+X)^{-1}$ therefore puts a neighborhood of the
identity in the closure of the rational group. The product of the
compact unitary groups is connected, hence the closure is the entire
group. The finite action on this coefficient module is trivial. Thus
an invariant rational subspace over $\mathbb C$ is invariant under
the real group by continuity and is zero or all of $M$ by its given
irreducibility. The elements of nonzero trace include the preimage
of an open identity neighborhood, since the trace at $1$ is
$\dim M>0$. For clarity, take the subgroup generated by this
preimage. Its image closure contains an open identity neighborhood
and hence the full connected group. Every rational element then
differs from an element of the subgroup by an element mapping into
the original neighborhood. This proves the required abstract
generation statement as well as density.
% END P74I1

% BEGIN P74I2
\paragraph{P74I2. One unit comparison on a fixed original record.}
Write $v_t$ for the stored theta datum, and $P\mapsto E_t(P)$ for
application of the creation polynomial to $v_t$ in the identity
source basis. At the base parameter $t_*=\mu_0$ this map is injective.
Indeed, the creation and annihilation operators on the Fock vacuum
satisfy the canonical commutators. In dual coordinates, applying
an annihilator to $E_{t_*}(P)$ applies the corresponding partial
derivative to $P$, up to a fixed nonzero normalization. The real
basis is orthonormal, so these coordinates give all derivatives.
A nonzero term of highest total degree can be differentiated to a
nonzero constant with its own multi-index. The image of $1$ is
nonzero. For the base finite datum take the rectangular identity on its
supported half. This is possible since $\operatorname{rank}W\le n$.
Its lower block vanishes and its top determinant is $1$ at every
residue prime. Both finite factors are nonzero. The vacuum has
unit norm, and the half-lattice Schrödinger realization on this
nonzero vacuum is injective (its distinct translates are separated
by lattice characters). Thus a polynomial in the kernel would have
all these constant derivatives in the kernel, an impossibility.

Fix an allowed identification of letters $\ell$ and the lowest-type
vector $u$ specified by the tie. Denote by $P_O(v)$ the complex
coefficient extension and renaming of $O.\mathrm{io}(v)$. On the
complex coefficient module $M$ form the two readings
\[
 A(w)=\mathrm{ptP}_{t_*}(\mathrm{ioP}(u,w)),\qquad
 B(O.c'_0 v)=E_{t_*}(P_O(v)).
\]
The second formula extends complex linearly. In fact the coefficient
tie identifies $O.c'_0$ on a basis extended from the lower field
with the original complex basis. The basis remains independent and
spanning on extending coefficients. The same argument applied to
the original rational unitary action shows that $O.c'_0$ intertwines
it with the action on $M$. This can be checked on the lower basis.

The extension and extremal-vector pins give a unit vector $w_0$ and
$\mathrm{ioP}(u,w_0)=\Delta$, the specified product of extremal
minors. The source basis in $A$ is the stored $B_V$.
Each $B_V$ is unitary as $V$ has identity Gram matrix. Take the
real group element $g$ with those components, and let $\chi$ be the
unit scalar in the real oscillator basis-change formula. Compact
equivariance, applied also to $g^{-1}$, gives
\[
 A(w_*)=E_{t_*}(\Delta),\qquad
 w_*=\rho(g^{-1})(\chi w_0),\qquad \|w_*\|=1 .
\]
Set $z_*=O.c'_0(O.\mathrm{hw})$. Its norm is one: the coefficient
pairing of the tensor of the pinned extremal vectors with itself in
the identity Gram matrix is $1$. The polynomial pin for $O.\mathrm{hw}$
is the product of the geometric minor polynomials. In the positive
embedding its row labels start at zero. In the conjugate embedding
they start after the $p$ positive rows. The weight tie gives the
successive highest-weight differences as their exponents. Under each
permitted renaming this gives the same $\Delta$, so
$B(z_*)=E_{t_*}(\Delta)\ne0$. Alternatively, nonvanishing follows
from injectivity of $O.\mathrm{io}$, of the letter renaming, and of
$E_{t_*}$, since $z_*\ne0$.

For $\gamma\in U(V)(F)$ let $g_\infty$ be its real part. The finite
part acts trivially on $M$. Compact equivariance gives
$A\rho(\gamma)=\omega_{t_*}(g_\infty,1) A$. For $B$, apply the
rational substitution pin to $P_O(v)$ and then the real oscillator
substitution formula. The remaining scalar is
$c_O(\gamma_E)\chi(g_\infty)^{-1}$, the pinned rational polynomial
factor divided by the real oscillator factor. The two identities on
the spanning set $O.c'_0(v)$ extend to $M$. P74I1 consequently gives
$A=a_OB$ and $a_Ow_*=z_*$. Both vectors have norm one, so $|a_O|=1$.
The latter equation is independent of the letter identification, and
$w_*\ne0$ makes its scalar unique. The vector $u$ is also fixed,
since each of its components is a subtype with prescribed underlying
vector. Finally strip off $E_{t_*}$ using its injectivity and the
source-basis substitution formula. This gives the polynomial identity
before applying any tuple vacuum, so applying $E_t$ gives
\[
 \mathrm{ptP}_t(\mathrm{ioP}(u,O.c'_0 v))
       =a_O E_t(P_O(v))\qquad\text{for every }t,v.
\]
This is $\mathrm{IoAt}(O,v_0,a_O)$. If no letter identification is
allowed the assertion is vacuous and $a_O=1$ works.
% END P74I2

% BEGIN P74I3
\paragraph{P74I3. A single coefficient phase over all extensions.}
Fix the original lower map $b=c_{0,L}(i)$. If no pinned record
satisfies coefficient base change for this $b$, use $a=1$.
Otherwise choose one such record $O_*$, allowing an arbitrary tuple
set and coefficient extension, and take $a=a_{O_*}$ from P74I2.
Let $O$ be any other record satisfying both pins. Fix an allowed
letter map and lowest-type input. Apply its reading and the anchor
reading at the common base parameter and their respective extremal
vectors. Their left sides agree: the tensor $v_{\rm ext}$ of the
extremal vectors over the lower field maps to $O.\mathrm{hw}$ and
$O_*.\mathrm{hw}$ under the two field extensions, since coefficient
extension commutes with tensors and extremal minors. Both complex
coefficient vectors are consequently $b(v_{\rm ext})$.
On the right, the variables are renamed with the same map.
Coefficient extension of each determinant polynomial is the
determinant of the same matrix of indeterminates, and products and
powers commute with extension. Thus both right sides use precisely
$E_{t_*}(\Delta)$. It is nonzero by P74I2. Cancellation gives
$a_O=a$, and P74I2 then proves the full reading on $O$. This shows
$|a|=1$ and proves \texttt{exists\_scalar\_core} with its order of
quantifiers across all tuple sets. No coefficient field containing
both extensions needs to be chosen besides their specified
embeddings in $\mathbb C$.
% END P74I3

% BEGIN P74I4
\paragraph{P74I4. The cover value on every marked analytic fiber.}
Fix a pinned record $O$ and any scalar $c$ satisfying its original
$\mathrm{XiConst}(O,c)$. For an arbitrary tuple $t$, denote its
higher model by $U_2$. Let $k$ be the adelic element and $x_k$ its
stipulated determinant root in the test defining $\mathrm{CovAt}$.
We claim that at every framed point $y_2$ equipped with the prescribed
analytically marked fiber,
\[
              \mathrm{ev}(c_\delta,y_2)=c\,\xi_t(x_k).
\]
The higher dictionary specified in $\mathrm{XiConst}$ gives a
component, its frame, and a rational $\gamma$ such that the lattice
and frame there are those of $\gamma k$ and the marking is its
adelic level marking. The core fixes the CM type, orthonormal
reference frame and common Hermitian form. Write $D_\gamma$ for the
rational unitary matrix for $\gamma$. It maps the test lattice for
$k$ bijectively onto the component lattice for $\gamma k$ by the
finite-lattice formula, preserves the Hermitian form, and its
inverse carries the component frame back to the test frame. Period
labels agree after the permutation given by the two ordering maps.
The two level markings agree after this change: at each completed
lattice their difference lies in the level multiple, by the two
adelic marking congruences. A rational vector in this multiple at
every completion belongs to the level multiple of the rational
period lattice. Thus the resulting polarized complex tori agree as
marked moduli objects. Their points in the higher fine moduli space
are equal. The section $c_\delta$ has weight and index zero, so
its evaluation depends only on this common underlying point.

In the value law at that dictionary point a rational unitary factor
$\gamma'$ may occur, from the representatives in $\mathrm{XiConst}$.
Its determinant has a root $b\in E^\times$ by the rational
determinant-root law (Hilbert 90). The product of its diagonal
image with $x_k$ is a root for $\gamma'k$. The tuple pin reads the
dictionary character as $\xi_t(bx_k)$, which equals $\xi_t(x_k)$
by multiplicativity and triviality on $E^\times$. Substituting
in the given $\mathrm{XiConst}$ value gives the asserted formula
with the original $c$, including $c=0$. This proves
\texttt{covAt\_of\_xiConst\_core} for the whole marked fiber.
% END P74I4

% BEGIN P74I5
\paragraph{P74I5. The unit attached to each component.}
Take a pinned record at the base parameter. Its embedding on the
lower field is the distinguished one, and its lower dictionary,
character $\mu$ and frame are the corresponding ones of $T$. Fix a
component with period lattice $L$. Choose a symplectic basis of $L$,
compatible coordinate labels and a classical vacuum $v$ for
$\mathcal O_E^n\otimes L$. In these coordinates $\Theta(v)$ denotes
its classical theta function. There are two comparisons to make.

First let $\vartheta_x$ be the analytic reading of the geometric
generator in a unit frame $x$. It has the canonical multiplier and
Gaussian-weighted squared mean one. The $n$-fold product and
$\Theta(v)$ belong to the same vacuum line on the product torus.
Indeed, the Heisenberg lattice acts on both by the prescribed
semicharacter, and the holomorphic space for a principal polarization
is one-dimensional, also for its external product. Thus
$\prod_b\vartheta_x(z_b)=C(x)\Theta(v)(x,z)$. The vacuum mean
formula gives $|C(x)|=\rho_L>0$ independently of $x$: integrate
the squared norms with the same Gaussian metric, using the product
probability measure on the torus. Equivalently one can compute the
classical vacuum norm in a symplectic basis by orthogonality of its
Fourier modes followed by the Gaussian integral. The normalized
vacuum norm at a unit frame is independent of that frame.

Here $C(x)$ can be made constant on the component. To see the phase,
translate by the level division points and take degree-zero theta
values at the origin. These give pairs of holomorphic readings
$(S_q,T_q)$ of the same weight, with $S_q=C(x)T_q$ on unit frames.
At every unit frame some $T_q$ is nonzero. More explicitly the
specified half-level lattice $\mathfrak m L$ is self-dual for
the original alternating form divided by $M$, by
$\mathfrak m\overline{\mathfrak m}=(M)$ and self-duality of $L$.
On that lattice the original theta function is of degree $M\ge3$.
The evaluations on its $M$-division grid have trivial common
kernel. For a principal polarization the translations by
$M$-torsion act irreducibly on the degree-$M$ theta space:
in a symplectic basis the diagonal half separates the
$M^g$ residue Fourier modes, and the other half permutes them
transitively. Evaluation at zero is a nonzero functional since the
$M$th power of a principal theta bundle, with any flat twist, is
generated there. For example, take a product of $M$ translates of
the principal section, with prescribed sum of translation parameters
to realize the twist. For $M\ge2$ the parameters can be chosen
off the finitely many proper divisors imposed by vanishing of a
factor at zero. The common kernel of all translates of the evaluation
is invariant and proper, hence zero. The marked-root congruences
identify this grid with the required values modulo periods. The period multipliers are nonzero. Taking the same nonvanishing
coordinate in each column gives the product assertion.

Where $T_q\ne0$, the ratio is holomorphic and has constant absolute
value. The common weight passes this last assertion from unit frames
to arbitrary frames by triangular change of basis. It follows from
the open mapping theorem that the ratio is locally constant.
The local ratios agree on overlaps since they all give $C(x)$ on
unit frames and the common weight cancels. The period component is
connected, so there is a single $C\ne0$, with $|C|=\rho_L$.
Changing the basis and ordering only makes a simultaneous coordinate
change in the Hermitian form, the Gaussian, and the arguments of
$\vartheta$, so the comparison applies to every generator reading.

Second, if the component lattice occurs adelically, choose a finite
$k_0$ moving the base lattice to it. The vector
$\omega_{t_*}(1,k_0)\mathrm{vac}$ satisfies the moved lattice and
Fock equations. In the complex period coordinates these are exactly
the same vacuum-line equations for $v$: the finite character is
the integral semicharacter, and the Fock equation identifies the
real restriction with a Gaussian times an entire theta function.
Consequently the moved vector reads $\kappa_0 v$, with
$\kappa_0\ne0$ by the unit norm of the adelic vacuum. Its norm fixes
the modulus without any choice of a theta constant. Indeed let
$H_x$ be the frame Hermitian form and write
$F_x(z)=|\vartheta_x(z)|^2\exp(-\pi H_x(z,z))$. It is periodic
and its torus mean is $1$. In the real period coordinates the
squared norm of the adelic reading has mean
\[
   \frac{|\kappa_0|^2}{|C|^2}\,
      \operatorname{mean}_{z}\prod_{b=1}^{n}F_x(z_b)
       =\frac{|\kappa_0|^2}{|C|^2}.
\]
For the adelic inner product this change of domain is legitimate:
the finite part is invariant in squared norm under an open lattice.
Unfolding its finite cosets gives a real period fundamental domain,
with its probability Haar measure. The period-coordinate map takes
Haar probability to Haar probability, and Fubini evaluates the
product. Since the adelic norm squared is $1$, $|\kappa_0|=|C|$.

Set $\kappa=\kappa_0/\mu(k_0)$. A second finite mover $k_1$ to the
same lattice changes the moved vacuum by the stabilizer character
$\mu(k_1k_0^{-1})$. The real part of a general mover acts by
moving the frame and its determinant character. Thus at $h$ the
vacuum reads $\kappa\mu(h)v$. The base twist is trivial, so the
lower character is the one tied to the original Weil action.
Left multiplication by a rational unitary element leaves $\mu(h)$
unchanged. Therefore $C_k=\kappa/C$ is a unit for the component
$k$, and the moved adelic vacuum, on arguments of finite part zero,
is $C_k\mu(h)$ times the original Gaussian times the product of the
actual generator readings. On a lattice reached by no adelic point,
any unit serves. Choosing one unit per component proves
$\mathrm{CcIs}$ with all its frame, basis and reading quantifiers.
% END P74I5

% BEGIN P74I6
\paragraph{P74I6. The signed contribution of one split class.}
The isometry between the level lattice pairing and the trace of the
Hermitian pairing supplies a sign $\delta\in\{1,-1\}$ before any
tuple or representative is chosen. Fix now the test data of
$\mathrm{ClsAt}$: a level and residue function $\phi$ of the stated
type, a split-supported representative $h$, an arbitrary target
mover $k$, and roots $Y(q)$ of the given higher marking.
Write $\mathcal Q=(\mathrm{LevQ})^n$, $\widehat\phi$ for the
unnormalized finite Fourier transform, $S=\mathcal O_E^n$,
$\Lambda$ for the prescribed source class lattice, $r=\operatorname{rank}W$,
and $y$ for the split determinant parameter of $h$. Let $v$ be
the classical vacuum in the test and suppose the original target
move $\omega_t(1,k)\mathrm{vac}$ reads $b v$.
Set $I=[\Lambda:\Lambda\cap S]$ and $\chi_W$ equal to the second
character of the tuple's original theta pair.

On the original half-lattice the finite theta test is the function
$\phi$ of the reduced matrix of columns. The lattice embedding
identifies rational columns in the standard model with the original
level lattice. Any finite lattice input is congruent to such columns
modulo $M$ by density in its completion. The level hypothesis puts
$M$ in all the residue primes used by the datum, so both tests have
the same reduced matrix and the same residue character values.
Outside its stipulated integral half the indicator vanishes. The
frame and its integral inverse identify the integral half with the
support of the corresponding lattice indicator (the primes of the
denominator are excluded by the level hypothesis). Fourier inversion
thus expresses the test as a sum of vacuum translates with weights
$\widehat\phi(q)/|\mathcal Q|$. To check the sign, at a rational
input $u$ the character of the translate by the idempotent half of
$\delta\iota(c_{\rm sep}\widetilde q)/M$ is the inverse-transform
character of the level pairing. Indeed the pairing of the input
with $c_{\rm sep}\widetilde q$ is minus $\delta$ times the trace
of $\xi_0$ times the Hermitian pairing taken in the opposite order.
Alternation cancels this minus sign in the Fourier character.
Summing the characters over $q$ then gives $|\mathcal Q|$ on the
matching residue and zero on the others.

Applying $\omega_t(h,k)$ transports each translation by its
symplectic action. The finite representative of $\delta Y(q)$
differs from this moved translation by twice a vector of the moved
integral vacuum lattice for every $q$ with $\widehat\phi(q)\ne0$.
Here is the congruence calculation. Let $e$ be the tuple's finite
idempotent, $X$ the move by $k$ of the full rational translate, $u$
the finite representative of $\delta Y(q)$, $L_1$ the lattice after
the target move, and $L'$ the lattice after both moves. The split
support of $h$ is disjoint from the selected residue primes, so it
fixes the $e$ and $\bar e$ components. The marking gives
$u-X\in L_1$ by multiplying the rational columns by $M$ and using
the common residues. Fourier support gives $(1-e)X\in L_1$: the
character of a nonzero Fourier coefficient is trivial on every
period of the given residue function. Also $MX\in L_1$. Consequently
$u-eX\in L_1$. Put $v'=(u-eX)/2$. The class lattice pins give
$N L_1\subset L'$ and $k_0L_1\subset L'$ for an auxiliary $k_0$
prime to $M$, while $M\mid N$. Thus $2k_0 v'\in L'$. The primary
root condition gives $M^j u=2Nw$ with $w\in L_1$ for some $j\ge1$.
It follows that $M^jX\in L_1$ and
$M^j v'=Nw-e(e/2)M^jX\in L'$, since the selected primes
exclude $2$ and $eL_1\subset L'$. As $M$ is odd,
$\gcd(2k_0,M^j)=1$, so B\'ezout gives $v'\in L'$.

For the central character, choose $s$ in the indicated power of the
half-level ideal which is a unit on the conjugate idempotent, so
$\bar e/s$ is integral and $su=2Nw'$ with $w'\in L_1$.
Such $s$ exists by the coprimality of the two halves. Since
$e\bar e=0$, one has
$\bar e v'=N\bar e(\bar e/s)w'$. In the symplectic pairing with
$eX$ only this conjugate component contributes. Transfer the factor
$M$ to the first entry: $eMX$ and
$(N/M)\bar e(\bar e/s)w'$ both lie in $L'$. The additive character
is therefore one by self-duality. As the half-pairing against
$2v'$ is the full pairing against $v'$, the Heisenberg product law
and the doubled-lattice vacuum equation identify the two
translations with no residual scalar.

For precision about the size and phase of the projection, first
consider a single split support. On the target tensor lattice take
representatives of the new lattice modulo its intersection with the
old. The finite sum of translates weighted by the new lattice
semicharacter, divided by their number $J$, is precisely the
new vacuum projection. Distinct representatives give orthogonal
translates of the unit old vacuum: self-duality separates their
difference by a character of the old lattice. Their unnormalized
sum consequently has squared norm $J$. In the split model the
Weil operator has the phase $\chi_W(y)^{-1}$ times a positive
multiple of that sum. This is the positive split intertwiner of
CW6: the two complementary halves give finite translate sums with
positive inverse square-root normalization, and averaging against
the new semicharacter leaves that same phase. Unitarity now gives
the positive multiple $J^{-1/2}$. Thus on the classical reading the
factor is $\chi_W(y)^{-1}\sqrt J$ times the normalized projection.
Both the old and new rational translates read their specified
classical Heisenberg translates: the finite representative of
the negative row vector of $z$ reads translation by $z$, with
the lattice semicharacter unchanged under negation.

For disjoint split supports the intermediate tensor lattice $B$
is the sum of $B\cap A$ and $B\cap C$, with $A,C$ the two
end lattices. This holds on each local factor since there one of
the two moves is the identity, and hence on the global lattices.
In the two successive finite averages the classes from $B\cap A$
act on the starting vacuum by its prescribed character and the
classes from $B\cap C$ are absorbed in the final average. All
resulting classes have the same number of preimages. Thus the
composition of normalized projections on the original vacuum
equals the final projection. The two relative-index factors
multiply, as is also seen by taking local indices at disjoint
supports. Induction gives the factor for the full support.
The relative index in the tensor lattice is $I^r$: locally the
target lattice is free of rank $r$, and indices multiply over
places. The determinants multiply, giving the one character
$\chi_W(y)^{-1}$ for $h$.

Pair this result with a homogeneous creation polynomial $P$ of degree
$d$. In Gaussian coordinates a creation operator on the reading is
the holomorphic direction operator with its Gaussian connection
term. At the origin the evaluation of the normalized creation
polynomial gives $(\sqrt{2\pi\xi_R})^{-d}$ times the classical
theta derivative. For a letter with direction $u$, differentiating the real
Heisenberg translation in directions $u$ and $iu$ gives on its
holomorphic factor the operator
\[
 (2\pi\xi_R)^{-1/2}\!
 \left[\tfrac12(D_u-iD_{iu})
            -\pi\sum_b H_x(u_b,z_b)\right].
\]
Here $D$ is the real directional derivative. Each $H_x(u_b,z_b)$
is antiholomorphic linear in $z_b$, so the holomorphic derivative
kills it, and its value at the origin is zero. Iterating on a
monomial and evaluating at the origin leaves only the classical
derivative, with one factor $(2\pi\xi_R)^{-1/2}$ per letter.
This is the same direction convention as the
classical harmonic reading. The finite translation of $\delta Y$
reads the classical translate by $-\delta Y$. Reflection fixes the
canonical even vacuum and commutes with the lattice projection.
Replacing all roots by their negatives replaces $z$ by $-z$ in the
classical theta function. A homogeneous order-$d$ derivative at zero
therefore acquires $(-1)^d$, so the effect of $-\delta$ is
$(-1)^d\delta^d$. By linearity the resulting value is
\[
 \frac{\delta^d(-1)^d}{(\sqrt{2\pi\xi_R})^d}\,
 b\,\chi_W(y)^{-1}\,\frac{\sqrt{I^r}}{|\mathcal Q|}
 \ \mathrm{thP}\!\left(
    \mathrm{gphi}(\Lambda,L,v,\widehat\phi,Y),P,x_k\right).
\]
Here $\mathrm{gphi}$ forms the indicated sum of classical
Heisenberg translates of the projected vacuum, and $x_k$ is the
moved frame. This is the one-class formula $\mathrm{ClsAt}$ on
the original lattices and characters.
% END P74I6

% BEGIN P74I7
\paragraph{P74I7. One point function and the identity level element.}
We now prove \texttt{point\_identity\_core}. If there is no pinned
record use the constant function $1$ and $z=1$. Otherwise choose
one pinned record for P74I5. Its lower model, reference frame and
dictionary come from $T$ at the distinguished embedding. In any
other pinned record $O$, restriction of its coefficient embedding to
the lower field is that same distinguished embedding. Its lower
dictionary and $\mu$ are transported along the given identification
$T.V=W$. Thus for each adelic $g$ its component unit agrees with
$C_{k(g)}$ of the anchor, where $k(g)$ is the component in the
anchor dictionary. This transport consists in substituting equal
embeddings and equal Hermitian spaces. The tables of higher
components are allowed to differ. Apply P74I3 for the fixed lower
coefficient map, and P74I6 for the tied lattice isometry. Put
\[
 d=\sum_{\tau}\sum_{j<\mathrm{mult}(\tau)}\mathrm{wtA}(\kappa_i)(\tau,j),
 \qquad K_d=a\delta^d(-1)^d(\sqrt{2\pi\xi_R})^{-d}.
\]
On $\mathrm{AutQ}(W)$, using the representative $q.\mathrm{out}$
specified by the point reader, set
\[
           \lambda(q)=K_d C_{k(q.\mathrm{out})},\qquad z=1.
\]
These are independent of $O,t,c$. The factors $a$ and $C_k$ have
modulus one, $\delta$ is an integer unit and $\xi_R>0$ is the
absolute value of the distinguished embedding of the nonzero
imaginary parameter. Hence $\lambda$ is everywhere nonzero.

Take $t$, placed coefficients $x$, an original section $f$ with
$\mathrm{IsCls}(O,t,x,f)$, and $g=q.\mathrm{out}$. The placement
pin supplies a definite form $y$ whose values are exactly the
coefficient readings $O.c'_0(x_j)$ on the prescribed representatives
$h_j$, and $O.\mathrm{thAd}$ is the reading of its theta lift.
Use the rational translate $\gamma g$ in the dictionary, a common
compatible pair of lower and higher markings, and its determinant
root $x_\gamma$. Self-duality of the period lattice gives the
dual roots $Y_j$ used by both short-series formulas. Write
$\kappa'$ for the factor taking the classical vacuum theta to the
product of the lower generator readings in this chart. Let $S(g)$
be the finite class sum in this common marking and frame, with
these $Y_j$, the stored Fourier coefficients, and $O.\mathrm{io}(x_j)$.

On the geometric side the marked lift and the section's class
value, as in T73C1--T73C3, evaluate the actual $f$ as
\[
       \mathrm{rdU}(f;O.\mathrm{dict}(g))
        =c\,\xi_t(x_\gamma)\,\kappa' S(g)
\]
by P74I4. On the adelic side unfold the compact definite quotient
into its specified double cosets. The weight of coset $j$ is the
given subgroup volume divided by its stabilizer size $\gamma_j$.
Indeed its quotient measure is that subgroup volume modulo the
finite stabilizer, and the opposite level characters of kernel and
form cancel on the subgroup. P74I3 supplies the factor $a$. P74I5
reads the base moved vacuum as
$C_{k(g)}O.\mu(g)\kappa' v$. The twist from $t_*$ to $t$ multiplies
this by $\xi_t(x_\gamma)$. This follows by twisting the Weil action
on the determinant root and using its uniqueness on theta functions. The base character is trivial. Apply P74I6 to each coset.

For clarity, if $m_j$ and $Z_j$ denote the lattice's stored
$\mathrm{mIdx}$ and $\mathrm{zCard}$ at $N_j$, its arithmetic
factor satisfies
\[
  \iota_O(\mathrm{ep}_j Z_j/m_j)
     =\chi_W(y_j)^{-1}\sqrt{[\Lambda_j:\Lambda_j\cap S]^r}.
\]
Here $S=\mathcal O_E^n$ and $y_j$ is the supported determinant
of $h_j$. The ideal of its supported finite idele gives the
class lattice's determinant up to conjugation. Outside its support
the characters are unramified or trivial on units, by the bad-prime
pin. The original tuple ideal-character pin gives its character
value, and the squared norm pin for this ideal supplies the positive
square root of the index (RTG9). Conjugating the class lattice
preserves its index against $S$. Substitution and the definition of
$\mathrm{ep}_j$ give the displayed equality, with the same
$Z_j,m_j$ as in the geometric class sum. The Fourier coefficient
of the finite test is the embedded stored weight by the tuple
Fourier pin (RTG11). Thus summing the formula of P74I6 gives
\[
 \mathrm{thF}_t(y)(g)=K_d C_{k(g)}O.\mu(g)\,
                         \xi_t(x_\gamma)\kappa' S(g).
\]
Multiplying by $c$ and substituting the geometric value proves the
identity even if $S(g)$ or $c$ vanishes. For $z=1$ the full
level-action laws give $\mathrm{actLev}(1)f=f$, so the product
$O.\mu(g)\mathrm{rdU}(f;O.\mathrm{dict}(g))$ is exactly
$\mathrm{rdQ}(O,1,f,q)$. For every tuple choose the root of unity
$1$ in the distinguished original field, of order one. We have
proved for all the data in the original quantifier order that
\[
      c\,O.\mathrm{thAd}(t,x,q)
        =\lambda(q)\,\mathrm{rdQ}(O,1,f,q).
\]
The section and the coefficient lift here are the ones given by the
test. Neither surjectivity of a space of sections nor spanning by
class sums is needed for this identity.
% END P74I7

\section*{A prescribed base residue and the descent from split short series}

Fix the base tuple $\mu _0$, its supplied record $O$ with the original
$\mathrm{Hk3.Core}$, tie and embedding pins, and
$K=O.c'.F$ with the specified embedding $\sigma:K\hookrightarrow
\mathbb C$. Write $n=\operatorname{rank}V_i$, $s=\operatorname{rank}W_i=D.R.r=r$, $b=O.\mathrm{bad}(())>0$,
$\mathcal O=\mathcal O_E$, and $\mathcal L=D_w.\mathrm{La}$.
We use the standard $V_i$ frame, the pinned $W_i$ lattice embedding
and $\mathrm{HalfIs}$ at this tuple. The finite datum in the spanning
argument has to do two jobs. Its
Fourier transform weights a geometric short series, and its untransformed
indicator has to be the precise finite oscillator datum on the base tuple.

% BEGIN S74B1
\paragraph{S74B1. A generating residue tuple with the prescribed finite indicator.}
Put $\mathfrak a=\prod_{P\in P_g}P$ and
$Q=\mathcal L/\bar{\mathfrak a}\mathcal L$.
The pinned half-level identifies $\bar{\mathfrak a}$ with
$M_{\rm fr}(l)$. There exists a tuple $x_1\in Q^n$ generating $Q$
over $\mathcal O$ for which the original finite test function at
$\mu _0$ is exactly $\mathbf1_{x_1}$ on residue tuples.

Indeed the base twisting character is trivial. Its ramification
criterion forces $Q(\mu _0)=\varnothing$, so its factor in the finite
weight is the empty product. Choose the supplied idempotent over
$P_g$. At $P\in P_g$, let $k_P=\mathcal O/P$ and let
$\rho_P:\mathcal L\to k_P^s$ be the reduction of the integral row
pairings in the $W_i$ frame. Its scalar rule is
$\rho_P(c x)=(\bar c\bmod P)\rho_P(x)$. We verify that the joint
map onto $\prod_P k_P^s$ has kernel $\bar{\mathfrak a}\mathcal L$.

Here are the denominator details for the lattice embedding $\iota$.
Let $L$ denote the self-dual frame lattice, $N=\operatorname{Nm}(\mathrm{ccK})$
and $m=\operatorname{Nm}(\mathfrak b)$ in the host data.
The defining lattice inclusions give
$N L\subset\iota(\mathcal L)$ and $m\iota(\mathcal L)\subset L$.
The half-level primes avoid $Nm$ by the coprimality with the pinned
bad factors and the half-different condition. Use finite dual
generators $s'_j$ of $L$. Their row pairing matrix $B$ has an integral
adelic left inverse $A$ with $AB=1_s$. More explicitly a finite
projective dual family resolves the identity on $L$; in frame
coordinates its evaluation rows give $B$ and its coefficient columns
give $A$. The coefficients $c_j$ of $y=\sum_j\bar c_j s'_j$ are the
entries of its pairing row times $A$. Choose $z_j\in\mathcal L$ with
$\iota z_j=Ns'_j$. For prescribed rows $t_P$, lift the conjugate
coefficients of a linear combination of the $z_j$ by the Chinese
remainder theorem, prescribing the row $N^{-1}t_PA_P$ at $P$.
The resulting pairing rows are $t_PA_PB_P=t_P$.

If all pairings of $x$ vanish modulo the indicated primes, the
coefficients $c_j$ of $m\iota x$ vanish at every $P$ by the same
coefficient formula. Thus $Nm\,x=\sum_j\bar c_j z_j$ is in
$\bar{\mathfrak a}\mathcal L$. Pick $v\in\mathcal O$ with
$Nm v\equiv1\pmod{\mathfrak a}$. The identity
$x=\overline{1-Nm v}\,x+\bar v(Nm\,x)$ gives the desired inclusion.
The reverse inclusion is the scalar rule. This proves the kernel and
surjectivity statements including the case of an empty set of primes.

Let $U_P$ be the reduced $V_i$ frame matrix. Both it and its inverse
are integral before reduction. Since $s\le n$, write
$J=\binom{1_s}{0}$. Lift the rows of $U_PJ$ simultaneously through
the joint map, and call their classes in $Q$ the entries of $x_1$.
For any pairing row $t_P$, the row $(t_P,0)U_P^{-1}$ gives
coefficients expressing $t_P$ in the lifted rows at $P$.
Conjugate Chinese remainder lifts make these coefficients simultaneous,
so the entries span $Q$. For an arbitrary tuple with matrix of rows
$T_P$, its residue equals $x_1$ exactly when $U_P^{-1}T_P=J$ at
every $P$. These are the upper identity and lower zero tests defining
the finite oscillator factor at the base. Both its values, zero and
one, are preserved under $\sigma$.
% END S74B1

% BEGIN S74B2
\paragraph{S74B2. A field of definition for the short-series move.}
Suppose a rational unitary move carries $(\Lambda,N,P,w)$ to
$(\Lambda',N',P',w')$, with the integral harmonic, level and
connection hypotheses of the original short-value construction.
Take $P'$ to be the pushforward of $P$ and the weights to agree on
moved residues. Write $U,U'$ for the two normalization-one short
values. They are natural functions of framed points by T73C1.
The analytic change-of-vacuum input on the pinned complex charts
takes the following precise form, with a single $|\tau|=1$:
\[
 U'(y)\sqrt{a}=\tau\sqrt{a'}\, U(y),\qquad
 a=[\Lambda:\Lambda\cap\mathrm{Std}]^{p+q},\quad
 a'=[\Lambda':\Lambda'\cap\mathrm{Std}]^{p+q}.
\]
It suffices to give this identity on analytic representatives for frame
changes since $P,P'$ have the same weight. Then there is one $z\in K$
such that for \emph{every} commutative $K$-algebra $S$ and framed
$S$-point $y$, $U(y)=zU'(y)$ (with $z$ mapped into $S$).

To see this, the indices are positive and hence the complex ratio is
$c=\sqrt{a}/(\tau\sqrt{a'})$. If $U'$ vanishes at all complex
points, both readings vanish there and we take $z=1$. Otherwise
choose a point where $U'$ is nonzero. For each $K$-automorphism
$\alpha$ of $\mathbb C$, transport the frame by $\alpha$. Naturality
and the identity at its transported point give
$(\alpha(c)-c)\alpha(U'(y))=0$. Thus $c\in\sigma(K)$. Here the
fixed field is $\sigma(K)$: a separable algebraic element outside it
has a different conjugate, and a transcendental element can be
included in a transcendence basis and translated by one. The resulting
isomorphisms extend to the algebraic closure. Take $z=\sigma^{-1}(c)$.

For the extension to test algebras, use the framed chart cover of
TDR3 and TDR5--TDR7; D74U2 supplies its section law at power one.
Over each maximal localization $S_{\mathfrak m}$ a chart containing
the image of the closed point contains the entire image. The universal
frame of its reduced finite type $K$-chart $A$ and the given frame
therefore pull back to a common frame via $A\to S_{\mathfrak m}$
and $S\to S_{\mathfrak m}$.
On $A$ the difference $U-zU'$ dies under every $A\to\mathbb C$.
Each maximal residue field is finite over $K$ and embeds in
$\mathbb C$ over $K$. The intersection of the maximal ideals is the
nilradical, zero in $A$. The difference is therefore zero in $A$.
Naturality gives equality in every maximal localization of $S$, hence
in $S$ itself. This explains why the argument permits nonreduced $S$.
The analytic vacuum identity is the remaining input to this descent.
% END S74B2

% BEGIN S74B3
\paragraph{S74B3. From the identity cover to a single generating residue.}
Assume the stable short-series span on the pinned identity cover, at
$n\ge 4[E:\mathbb Q]r+4$, and the complex vacuum identity in S74B2 for
rational unitary transport. Every cusp reading then lies in the complex span of raw
sections with the \emph{one} weight $w=\widehat{\mathbf1}_{x_1}$ of
S74B1. Each such section uses a unimodular $\Lambda$ trivial at the
bad modulus and at $l.M$, a permitted $N$, and
$P\in\mathrm{harmL}(K,\Lambda)$ of the given weight.

In fact an identity-cover generator with lifted tuple $\tilde x_0$
has normalized short value $f$ with scalar $\beta_H(\Lambda)\ne0$.
Its multiplier is one on the identity cover. Evaluating there on
geometric fibers and using uniqueness shows that $f$ has the required
normalized value on every reduced test chart (geometric points
separate elements of that chart). Multiplying the defining short
identity by $\beta_H$ gives the raw normalization-one value for
$\beta_H f$. The Fourier sum against the point indicator reduces to
its single summand, the level character of $\tilde x_0$; this
character depends only on its residue $x_0$. Admissibility of the
lifted tuple says precisely that its entries generate after passage
to the quotient. Hence allowing these $x_0$ already spans all the
cusp readings in question. The numerical bound in the supplied
record is exactly this bound since $[E:\mathbb Q]=2f$.

To construct the move, work first on the residue module of S74B1.
On each factor both tuples give surjections from an $n$-dimensional
space to the same $s$-dimensional space. Choose bases of their
kernels and lifts of a common basis of the target. The change of
these bases carries one surjection to the other. In a split pair the
unitary matrix is specified by its matrix on one factor, with inverse
transpose on the conjugate factor. Prescribe the first factor so the
adjoint sends the tuple $x_1$ to $x_0$; the two prescriptions are
compatible by construction. Lift the invertible residue matrices to
integral invertible local matrices and choose identity matrices at
the other places over $2b$. Weak approximation in the rational
unitary group gives a rational unitary matrix in these finitely many
open congruence neighborhoods (use the rational Cayley chart after
perturbing the local matrices away from its excluded divisor).
Its action is integral with integral inverse at these places. For
transported residues $q,q'$ the level pairings satisfy
$\langle x_1,q'\rangle=\langle x_0,q\rangle\pmod M$ by the
adjoint identity. Exponentiating proves precisely equality of the
Fourier weights. The resulting matrix preserves the integral
harmonic module by pushforward, and
triviality at $2$ and at the level supplies an allowable $N'$.
If $s_0$ is the old raw section, S74B2 gives $s_0=zU'$ on every test
chart. For $z=0$ its reading is zero. Otherwise $z^{-1}s_0$ is still
a section of the same weight, and the displayed identity says at
every test point that it is a raw section for the new tuple. Thus its
$K$-span contains $s_0$. This does not require constructing a
section from an arbitrary function on complex points. T73C2--T73C3
supply the original pinned short rules used on the identity cover.
The stable analytic span used here still needs the reproducing
kernel and level-to-short-series comparisons specified in ICL14 and
ICL16.
% END S74B3

% BEGIN S74B4
\paragraph{S74B4. Moving a raw lattice to a split finite point without moving its weight.}
Let $b$ be the supplied bad integer, $M=l.M$, $M\mid b$, and
$2\mid b$. If $\Lambda$ is any of the raw lattices in S74B3,
there are a pure finite $h$ supported on split places away from $bM$
and a rational unitary $d$ such that
\[
       d\Lambda=\overline{L_h},\qquad
       L_h=\mathrm{hLat}(V_i,h),\qquad d\equiv1\pmod{bM}
\]
in the columnwise congruence sense (denominators coprime to $bM$).
In particular the Fourier weight $\widehat{\mathbf1}_{x_1}$ can be
used on both sides of this move.

For clarity construct first a finite adelic unitary $g$ carrying the
standard completed lattice to $\widehat{\bar\Lambda}$ and equal to
one at primes over $bM$. Conjugation preserves the standard lattice,
unimodularity and level triviality. The two lattices coincide at the
places over $bM$ by triviality and at all but finitely many places by
commensurability. At the remaining places use an integral orthonormal
frame of the self-dual lattice. Here the host gives an unramified CM
extension at finite places. At a split place use a basis in one factor
and its Hermitian dual in the other. At an unramified quadratic field
place the residual nondegenerate Hermitian form contains a vector of
nonzero norm: if all norms vanished, polarization with all residue
scalars and nondegeneracy of the residue trace pairing would kill the
entire form. Lift this vector. Its norm is a unit, and the norm on
units is surjective by the finite-field norm followed by successive
correction on principal units using the residue trace. Rescale to norm
one, split off the orthogonal line, and induct. This gives the required
local frames and hence $g$.

Apply the finite factorization (BTH.49) at modulus $bM$, writing
$g=q_f h k$ with $q$ rational unitary and $k$ in the standard
lattice stabilizer congruent to one at the modulus. The split support
of $h$ gives $h=1$ there. Consequently $q\equiv1\pmod{bM}$
there, hence also as a rational matrix after clearing denominators
coprime to $bM$. Since $k$ preserves the completed standard lattice,
intersecting the lattice identity with the rational space gives
$\bar\Lambda=qL_h$. Unitarity gives $\bar q^{-1}=q^{\mathsf T}$,
so $d=q^{\mathsf T}$ has precisely the claimed direction.
The residue action of $d$ modulo $M$ is the identity, so any function
of the residue, including this Fourier weight, is preserved.

The split-support lattice argument of RTG7 supplies a permitted $N'$
and the original lattice guard for $L_h$; the conjugation in the last
display agrees with the lattice in the short value. The change of
variables pushes $P$ to its integral harmonic polynomial on
$d\Lambda$: pullback on linear letters commutes with contractions by
the Hermitian form and identifies the two coefficient lattices.
Under the analytic change-of-vacuum input of S74B2, the resulting
short value $t$ on the original dictionary thus satisfies
$L s_0=\sigma(z)t$ for the original raw section $s_0$,
where $L$ is its reading followed by evaluation on that dictionary.
% END S74B4

% BEGIN S74B5
\paragraph{S74B5. Reduction of the base class-sum span to the original analytic comparisons.}
Keep the full pinned base record. Let $\mathcal S$ be its cusp
sections for which some coefficient tuple satisfies both the original
place condition and the original geometric class-sum rule. Assume
the stable span and transport inputs of S74B3 and the two point
comparisons below. Then every cusp section lies in
$\operatorname{span}_K\mathcal S$, as required by
$\mathrm{ClassSumsSpan}$.

Write $\ell=\mathrm{lieFac}>0$ and
$a(g)=c_{\mathrm{dict}(g)}\ell\,\mu(g)$, using the original frame
constants and unitary character. For a moved lattice $\bar L_h$,
let $t(g)$ be its normalization-one short value. The harmonic pin
expresses its polynomial as $O.\mathrm{io}(v)$ with $v$ in the
prescribed coefficient space. The needed one-class kernel comparison
is
\[
     K_{h,v}(g)=A_h\,a(g)t(g),\qquad
 A_h=\mathrm{cmpConst}(D,k_a,\delta,\deg W)\,
         \ell^{-1}\sigma(\mathrm{epF}(L_h)).
\]
The constants can be divided out: $k_a$ has norm one, $\delta$ is a
sign, and the lattice guard and the split determinant ideal
$\mathfrak c$ give
$\mathrm{epF}(L_h)=B(\mathfrak c)^{-1}q$ with $q>0$ rational.
The ideal is nonzero and coprime to the bad modulus; the pinned value
of $\sigma(B(\mathfrak c))$ is a unitary value times a positive
square root of its norm. In addition $a(g)\ne0$ for each $g$.

By K73G3, $K_{h,v}=\sum_j b_j\Theta_j$ for finitely many lifts of
actual rational form inputs. For each input choose its base-class
coordinates by the pinned representative presentation. This witnesses
the original place condition. The geometric class-sum construction of T73C5
furnishes a section $f_j\in\mathcal S$ on these coordinates.
The second needed point comparison, in the \emph{same} frame
normalization as the first, is $\Theta_j(g)=u\,a(g)Lf_j(g)$,
where $u$ is the common norm-one phase for the base record. Substitution
and pointwise cancellation of $a(g)$ give
\[
       t=\frac{u}{A_h}\sum_j b_j Lf_j .
\]
This calculation allows zeros of $t$ and zeros of individual lifts.
S74B4 places every raw reading in the same span, and S74B3 places
every cusp reading there. Finally K73G4 descends this finite span
through $L$ to the coefficient field $K$ at the original cusp
sections.

This is a reduction of the base spanning requirement. The stable
reproducing kernel with its chart comparisons, the analytic
change-of-vacuum identity, and the geometric comparisons on the same
pinned normalization are substantive inputs. Finite-point generation
and coefficient descent alone give no identity for those inputs.
% END S74B5

\section{From split point kernels to the remaining spanning problem}

% BEGIN J74T1
\paragraph{J74T1. The one-class normalization for an arbitrary split move.}
Fix the original pinned base record and a finite split move $h$ off its
bad integer as in S74B4. Write $L_h$ for its unconjugated lattice,
$S=\mathcal O_E^n$, $I=[L_h:L_h\cap S]$ and $r=D.R.r$.
For the geometric lattice $\Lambda=\bar L_h$ let $m$ and $Z$
denote $\mathrm{mIdx}(\Lambda,N,r)$ and
$\mathrm{zCard}(\Lambda,N,r)$ at an allowed denominator $N$.
Let $\chi(y)$ be the original lift character on the supported
determinant root $y$. Then
\[
        \sigma\bigl(\mathrm{epF}(L_h)Z/m\bigr)
                       =\chi(y)^{-1}\sqrt{I^r}.
\]
Here $\sigma$ embeds the coefficient field of the base record. The
identity applies to a move produced after the raw short series was
chosen; $h$ need not be one of the representatives originally stored
in the class sum.

Indeed, put $d_P=v_P(\det h_f)$. Unitarity gives
$d_{\bar P}=-d_P$ and off the moved split pairs $d_P=0$. Define
$\mathfrak c$ by the exponents $\max(d_P,0)$. The support and
coprimality hypotheses give
$\mathfrak c+\bar{\mathfrak c}=\mathcal O_E$ and
$\det(L_h)\bar{\mathfrak c}=\mathfrak c$. At a moved pair the
unramified local character values on uniformizers are inverse since
the character restricts trivially to the diagonal split base field
there. A negative supported determinant valuation therefore gives the
positive exponent at the conjugate prime in $\mathfrak c$; a positive
one gives that exponent at the original prime. At the other places
the supported determinant is one. This proves that the ideal character
on $\mathfrak c$ has value $\chi(y)$. By the two pins of RTG8--RTG9,
for $Q=\mathrm N(\mathfrak c)$ there is a \emph{positive} rational
$q$ with
\[
 q^2 I^r=Q,\qquad
 \sigma(\mathrm{epF}(L_h))=q/(\chi(y)\sqrt Q).
\]
The level guard puts $NS\subseteq\Lambda\cap S\subseteq\Lambda$.
Multiplicativity of indices in this chain gives $Z=mI^r$:
conjugation preserves $S$ and the relevant index. All indices are
positive and finite since the lattices are full and $N\ne0$.
Thus the left side equals $\chi(y)^{-1}q I^r/\sqrt Q$,
which is the claimed expression since $\sqrt Q=q\sqrt{I^r}$.
% END J74T1

% BEGIN J74T2
\paragraph{J74T2. The two point comparisons use one common multiplier.}
Work at a pinned base record with the core, frame and oscillator
extension hypotheses of P74I2--P74I7. Retain the original prime
agreement, family laws, and inclusion of the embedded coefficient
field in $O.c'.F$ required by T73C5 for the chosen record.
Take a normalization-one
short value $t(g)$ after the move of S74B4, at the dictionary frame
of an arbitrary target adelic point $g$. The harmonic pin represents
its polynomial by $P=O.\mathrm{io}(v)$.
Use the indicator $\phi$ of S74B1. We can now supply both point
comparisons required by S74B5 on these data.

To see the normalization geometrically, use the marked period chart
of the given frame. Write $\mathcal Q=(\mathrm{LevQ})^n$,
$\kappa'$ for the generator-to-vacuum comparison in this frame,
and $H$ for the evaluation of the classical derivative polynomial
on the projected vacuum translated by the dual roots and weighted
by $\widehat\phi$. In the complex model the normalization-one
short value is
\[
            t(g)=\frac{Z}{m|\mathcal Q|}\,\kappa' H .
\]
Here is how its three finite factors arise. Use representatives for
the finite isotropic lattice quotient, using the same even lifts
and the same marked roots as in P74I6 on the
complex fiber. They provide an admissible presentation in the
defining short-value rule; uniqueness in T73C1 permits this choice.
The canonical torsion lifts read as Heisenberg translates of the
theta generator. In the pulled-back chart the scaled variable is
$Nz$, so the homogeneous derivative of order $d$ acquires $N^d$
by the chain rule. This cancels its defining $N^{-d}$ factor.
The harmonic projection removes the quadratic frame correction under
differentiation, as in T73C4. The derivative is linear in the section,
so it commutes with both
finite sums. The sum of translated vacua over the isotropic quotient
is $Z$ times the averaging idempotent applied to the vacuum. The residue sum gives the
weighted translate of that projection, with the factor
$1/|\mathcal Q|$. The remaining normalization in the geometric
rule is $1/m$, and reading the product generator supplies $\kappa'$.
These identifications use the chart's actual basis, marked roots,
theta connection and renamed polynomial; changing the marking in
just one of the two sums would change $H$.

Write $C_{k(g)}$ for the component constant of P74I5 and
$K_d=a\delta^d(-1)^d(\sqrt{2\pi\xi_R})^{-d}$ as in P74I7,
now using the coefficient phase $a$ of this one base record.
The rational element moving $g$ into the marked chart preserves the
adelic theta evaluation. The base twist is trivial. P74I2 and P74I6
give for the point kernel
\[
 K_{h,v}(g)=K_d C_{k(g)}\mu(g)\,
       \chi(y)^{-1}\sqrt{I^r}\,\kappa'H/|\mathcal Q|.
\]
The residue Fourier weight is the embedded geometric weight by RTG11.
Substitution of J74T1 gives
\[
 K_{h,v}(g)=K_d\, C_{k(g)}\mu(g)\,
                   \sigma(\mathrm{epF}(L_h))\, t(g).
\]
No geometric value has been divided out.

For an integrated lift $\Theta_j$ of a rational source vector take
its placed class coefficients and the section $f_j$ of T73C5.
Use this \emph{same} choice of $a$, sign, component constants and
charts in the finite sum of P74I7. Its calculation on a marked
adelic point precedes passage to the quotient representative, so it
applies here at each $g$. The pinned cover constant $c_0$ of this
base record is a root of unity. It gives
\[
 c_0\Theta_j(g)=K_d C_{k(g)}\mu(g)\,
                         \operatorname{rd}_{\mathrm{dict}} f_j(g).
\]
In both equations $C_{k(g)}\mu(g)$ is the same nonzero function.
In the notation of S74B5, $\ell=(\sqrt{2\pi\xi_R})^{-d}$,
so the required scalar for the integrated lifts is
$u=K_d/(c_0\ell)$, of modulus one. If
$K_{h,v}=\sum b_j\Theta_j$ as in K73G3, the equations yield the
somewhat simpler identity
\[
       t=\frac{1}{c_0\,\sigma(\mathrm{epF}(L_h))}
                 \sum_j b_j\operatorname{rd}_{\mathrm{dict}} f_j .
\]
The denominator is nonzero by J74T1. This proves generation for
every short value reached by a split move, including a value which
vanishes identically.

Consequently the stable identity-cover expansion and the analytic
change-of-vacuum formula retained in S74B2--S74B4 suffice for the
original $\mathrm{ClassSumsSpan}$ on this pinned base record: those
arguments bring the cusp readings into the span of the moved short
values, and S74B5 descends that span to its coefficient field.
The higher covers on the prescribed field are available under D74U5,
including the section for the trivial determinant character. T73C5
supplies class sums also after the allowed change of phase.
P74I7 supplies the common point identity. RTG16--RTG18 can therefore
transfer the resulting base span to the original all-tuple comparison
under their stated base, phase and coefficient hypotheses. The
analytic expansion and change-of-vacuum formula, and the outer
construction of one family with the required inputs, remain separate
existence requirements.
% END J74T2

\section{Changing a raw lattice on the original base dictionary}

% BEGIN V74M1
\paragraph{V74M1. Simultaneous roots for a rational column move.}
\label{claim:V74M1}
Let $M=l.M$, $\mathfrak b=l.Mb$,
$\mathfrak b\bar{\mathfrak b}=(M)$ and
$\mathfrak b+\bar{\mathfrak b}=\mathcal O_E$. Let $L$ be a marked
target period lattice with marking $m:D_w.\mathrm{La}\to L$.
Let $d$ be a rational unitary column change integral with integral
inverse over $M$. Let $N,N'$ be the two positive allowed short
levels, and let $\pi$ be the induced bijection on
$\mathcal Q=(D_w.\mathrm{La}/\mathfrak bD_w.\mathrm{La})^n$.
In applying the complex short-value formula one may choose the
original roots of both markings so that
\[
                  Y'(\pi q)=dY(q)\qquad(q\in\mathcal Q)
\]
as rational matrices, including when $N\ne N'$.

\begin{proof}
Represent the entries of $q$ by
$a\in(\bar{\mathfrak b}D_w.\mathrm{La})^n$, possible by
comaximality. At $P\mid\mathfrak b$ prescribe the coset
$Y(q)\in M^{-1}m(a)+L_P^n$.
At each $P\nmid\mathfrak b$ require
\[
       Y(q)\in 2N L_P^n\cap d^{-1}(2N'L_P^n).
\]
These local conditions equal the integral lattice away from
finitely many places. They admit a rational solution. Indeed a
single common denominator puts all cosets in one fractional
lattice, and each condition then specifies a congruence in that
lattice modulo a power of a prime. The Chinese remainder theorem
on its finite quotient solves the finite list simultaneously.

At $P\nmid\mathfrak b$ the vector $m(a)/M$ is integral: at
$P\mid\bar{\mathfrak b}$ its numerator is divisible by the full
$P$-part of $M$. Thus $Y(q)-m(a)/M\in L^n$. All places outside
$\mathfrak b$ already satisfy $Y(q)\in 2N L_P^n$;
increasing $j$ at the finitely many remaining places gives
$\mathfrak b^jY(q)\subset 2N L^n$. These are exactly the two
root conditions. Choose a half representative $a'$ of $\pi q$.
At each $P\mid\mathfrak b$, $a'-da\in M(D_w.\mathrm{La})_P^n$
since the localized half modulus there is $(M)$. Linearity of $m$
after localization gives $dY(q)-m(a')/M\in L_P^n$. Away from
this half the two terms are integral, and $dY(q)\in 2N'L_P^n$
by its construction. Increasing $j$ again supplies the second
root condition with $N'$. Invertibility of the residue action
makes $Y'$ a definition for every $q'$. The point formula permits
every admissible system of roots, so using this pair introduces
no new invariance assumption.
\end{proof}
% END V74M1

% BEGIN V74M2
\paragraph{V74M2. One unit for the projected vacua seen by the base dictionary.}
\label{claim:V74M2}
Use the genuine two-group Weil action of the pinned base tuple.
Let $S=\mathcal O_E^n$, $L_*$ be the target lattice of the base
vacuum $v_*$, and $v_k=\omega(1,k)v_*$ for a finite target mover.
Write $L_k=kL_*$ at the finite completions. For any of the self-dual
source lattices $\Gamma$ under consideration define $p_{\Gamma,k}$
to be the orthogonal projection onto the prescribed finite vacuum
character of $\widehat\Gamma\otimes\widehat L_k$; put
$u_{\Gamma,k}=p_{\Gamma,k}v_k/\|p_{\Gamma,k}v_k\|$.
For a rational unitary column move $d$, one scalar of modulus one,
independent of the target mover, intertwines the normalized vacua
for $\Gamma$ and $d\Gamma$ on all these points.

\begin{proof}
The projector is the probability Haar average of the finite
Heisenberg operators multiplied by the prescribed semicharacter.
The latter is the finite additive character of one half the
alternating pairing of $y$ with $w_0 y$, with the pinned sign.
Its restriction to rational columns is the factor
$\exp(\pi i\sum_b E_F(y_b,w_0 y_b))$ in the classical projector;
this follows from the product formula for the additive character.
The self-dual integral lattice makes the corrected operators a
group on that lattice. On a smooth vector the integral is a finite
orbit average. Density of the rational tensor lattice in its finite
completion identifies it with the original finite sum, with its
original denominator. In particular this comparison commutes with
reading the adelic vacuum in the classical model of P74I5.

For $k=1$ let $A=\widehat\Gamma\otimes\widehat L_*$ and
$B=\widehat S\otimes\widehat L_*$. Representatives of
$A/(A\cap B)$ give the summands of $p_{\Gamma,1}v_*$. They
have norm one. If two representatives differ outside $B$, an
element of $B$ separates their Heisenberg characters by self-duality;
unitarity then makes the two vectors orthogonal. Representatives
differing by $A\cap B$ give the same corrected summand. Hence the
squared norm of the average is $[A:A\cap B]^{-1}$, which is
positive. Writing $r=\operatorname{rank} L_*$, this index equals
$[\Gamma:\Gamma\cap S]^r$: locally $L_*$ is free of rank $r$
and the local indices multiply.

At the base, the simultaneous eigenspace for this finite vacuum
and the archimedean Gaussian equations is a line. In the
Schr\"odinger realization the two halves of the self-dual
Heisenberg lattice impose the support and translation equations:
the diagonal half restricts the function to one dual coset and
the translation half determines all its values there from one
value. One may choose integral symplectic coordinates at each
finite prime; self-duality supplies the dual pair and the fixed
semicharacter supplies the coset and the translation factors.
At infinity the Gaussian equations determine a scalar multiple
of the normalized Gaussian. Equivalently this is the
one-dimensional theta space for the principal vacuum polarization
used in P74I5. This argument applies to $A$ and to its image by $d$.

Use the rational source element $a=(d^{-1})^{\mathsf T}$; the
unitary relation gives $a=\bar d$. The native action is right
multiplication by $a^{-1}$ on maps from $V$ to $W$, hence it is
left multiplication by $d$ on their transposed matrix of rows.
It carries $A$ to the lattice for $d\Gamma$, preserves the
semicharacter by $E$-linearity and isometry, and preserves the
Gaussian equations since the source unitary group is compact at
infinity. Thus the unit vectors $\omega(a,1)u_{\Gamma,1}$ and
$u_{d\Gamma,1}$ are proportional. Define the unit $\tau$ by
$u_{d\Gamma,1}=\tau\,\omega(a,1)u_{\Gamma,1}$.

The finite target action also preserves the very same semicharacter.
Conjugating the probability average by $\omega(1,k)$ therefore
carries $p_{\Gamma,1}$ to $p_{\Gamma,k}$, including its
normalization. It follows that
$u_{\Gamma,k}=\omega(1,k)u_{\Gamma,1}$ for both source
lattices. The source and target actions commute because the
original action is a representation of their direct product.
Consequently
\[
          u_{d\Gamma,k}=\tau\,\omega(a,1)u_{\Gamma,k}
\]
with the one unit just chosen. The scalar comparing the classical
standard vacuum with $v_k$ in P74I5 multiplies both classical
unit projections by the same phase. It cancels, even if it
varies between components.
\end{proof}
% END V74M2

% BEGIN V74M3
\paragraph{V74M3. The raw short-value change at the required dictionary points.}
\label{claim:V74M3}
For the original base dictionary and the rational moves of S74B3
and S74B4, let $U,U'$ be the normalization-one short readings with
corresponding Fourier weights and with $P'$ the column pushforward
of $P$. On all those dictionary points the same $\tau$ of V74M2
satisfies
\[
               U'\sqrt{j}=\tau\sqrt{j'}\,U,\qquad
 j=[\Gamma:\Gamma\cap S]^r,\qquad
 j'=[d\Gamma:d\Gamma\cap S]^r .
\]
Here the levels $N,N'$ may differ, the two geometric raw sections
may have different Fourier weights, and zero short values are
allowed.

\begin{proof}
On a marked analytic fiber take the same period frame $x$ and
the same product-to-classical generator constant $\kappa_x$ for
both sides as in P74I5. In the next display $u_{\Gamma,k}$ denotes
the corresponding classical normalized projection. To pass from the
adelic reading to this convention, divide by the factor comparing
$v_k$ to the classical unit standard vacuum of this chart. Projection
is linear and its norm on each side changes by the same factor
$1/\sqrt j$, so this comparison factor is common to both source
lattices. Let $Z_N$ denote the finite isotropic
subgroup in the short-value sum, and let $m_N$ be its normalization
index $\mathrm{mIdx}(\Gamma,N,r)$. The short reading of T73C1
in this marking is
\[
 U=\frac{\kappa_x\sqrt{j}}{|\mathcal Q|}
      \sum_{q\in\mathcal Q} w(q)\,
          \mathrm{thP}_x(T_{Y(q)}u_{\Gamma,k},P).
\]
For the index factor, the original geometric sum uses denominator
$m_N|\mathcal Q|$. Reading the translated theta factors and summing
over $Z_N$ gives its cardinality times the vacuum projection: each
term is the weighted Heisenberg translate for its isotropic point
and all fibers of the finite averaging map have equal size. In
integral tensor coordinates $|Z_N|=m_N j$, by the index tower
through $\Gamma\cap S$ and local freeness of the rank-$r$ target
lattice. The norm of that projection is $1/\sqrt j$ by V74M2.
The three factors give the displayed coefficient. The connection in
the short sum contributes exactly the classical harmonic derivative
$\mathrm{thP}_x$: in normal coordinates Taylor expansion of the
translated product applies the indicated directional derivatives;
terms involving derivatives of its Gaussian gauge contract two
letters and vanish on the harmonic polynomial. This is the same
short-value reading used for a class in P74I6--P74I7, here applied
before any assumption on the split support of a representative.

Use V74M1 for $Y,Y'$. At infinity let $d z$ denote the column
change on each Lie-type block, using that block's embedding of $E$.
In the adelic theta model a rational pair acts by inverse
substitution on the argument. Indeed evaluation at zero is rationally
invariant, and evaluating a translated vector at zero evaluates
the vector at its translating argument. Heisenberg covariance
transports that translating argument by the symplectic action.
Thus for the finite translate and its transported counterpart
V74M2 gives equality of the classical functions at $z$ and $dz$
up to $\tau$. The Gaussian gauge in this comparison is unchanged
because $d$ is unitary. The finite Heisenberg translations match
literally by $Y'(\pi q)=dY(q)$, with the same sign convention on
both sides. The chain rule identifies differentiation at the origin
by $P$ on the composed function with differentiation by $P'$ on
the moved function: a direction in column $b$ is sent to the sum
of directions in columns $a$ with coefficients $d_{ab}$ in
the corresponding embedding. Apply this first to a monomial in
directions, then use linearity. Thus each new summand is $\tau$
times its old partner. The weight identity
$w'(\pi q)=w(q)$ and the residue permutation finish the calculation.

For the spanning application work with the complex reading map
$L=\mathrm{atDict}\circ\mathrm{rdU}$ throughout. Starting from
the stable span by identity-cover short sections with varying
generating residues, apply the preceding formula first to the
residue move and then to the split move. Each ratio is constant
on the entire original dictionary. The resulting split readings
therefore span every $Lf$ occurring in that stable span.
The coefficient-field descent for the final class sums is K73G4.
This use requires only the target lattices realized by the original
dictionary, all of which have a target mover by its marking
condition. S74B2 also describes descent of a uniform comparison on
all complex framed points, when that stronger input is available.
\end{proof}
% END V74M3

\section{A uniform overlap bound for the original finite level system}

% BEGIN F74N1
\paragraph{F74N1. Counting the self-dual lattices at each intersection index.}
Let $B$ be the rational alternating form of the given period chart
on $E^r$, of rational dimension $d=[E:\mathbb Q]r$. Its symplectic
basis makes $B$ nondegenerate. Fix one of its self-dual full
$\mathcal O_E$-lattices $L$. Write
$m(L,L')=[L:L\cap L']$. Among all self-dual period lattices there
are at most $m^{2d}$ choices of $L'$ with $m(L,L')=m$ for $m\ge1$.
The only lattice at index one is $L$.

Indeed, forgetting the $\mathcal O_E$-action gives an injection into
the set of full integral lattices self-dual for $B$. The integral
dual condition is the same condition on the underlying additive
group. Alternatively, finite generation and self-duality imply
fullness: the annihilator of the rational span would otherwise contain
a rational line lying in the integral dual and hence in a finitely
generated free abelian group. For a free abelian group of rank $d$,
there are at most $m^d$ subgroups of index $m$. To see this directly,
project onto the last $d-1$ coordinates. If the intersection with
the first coordinate axis is $a\mathbb Z$, the image is a lattice
$H$ of index $m/a$ in $\mathbb Z^{d-1}$. A lift of $H$ is
specified by a homomorphism $H\to\mathbb Z/a$, of which there
are $a^{d-1}$. By induction there are at most
$\sum_{a\mid m}a^{d-1}(m/a)^{d-1}\le m^d$ choices.
The rank-zero case is immediate.

Now choose $M=L\cap L'$, of index $m$ in $L$, and let $M^\vee$
be its integral dual. The inclusion $M\subset L$ gives
$[M^\vee:L]=m$ by transposing an integral basis matrix. Using
$M\subset L'$ gives the same index reversal
$[M^\vee:L']=[L':M]$. Since $[M^\vee:M]=m^2$, both of these
last indices equal $m$. For a given $M$, every possible $L'$ is
therefore an index-$m$ subgroup of $M^\vee$. Applying the preceding
count twice proves the bound; at index one the inclusions force
$L'=L$. The same argument works for any subset of the self-dual
lattices, in particular the original marked level set.
% END F74N1

% BEGIN F74N2
\paragraph{F74N2. The two-lattice inner product in the original finite level.}
Take the source level $D_\ell$ with half ideal $\mathfrak m$,
and finite matrix subgroup $J$. The level assumptions
are $(\mathfrak m,\bar{\mathfrak m})=(\mathfrak m,2)=1$.
Use the actual even column count $n\ge r$, $M>0$ lying in
$\mathfrak m$, a permitted source lattice $\Lambda\in V_{n,M}$,
and the specified normalized input vacua. For the finite test functions
$\Phi_L=D_\ell.\mathrm{phi}(L,\Lambda)$,
\[
 \langle\Phi_L,\Phi_L\rangle=|J|,\qquad
 |\langle\Phi_L,\Phi_{L'}\rangle|
       \le |J|\,[L:L\cap L']^{-n/2}.
\]

Here is the role of the marking in this estimate. Set
$S_L=\mathfrak m\bar{\mathfrak m}^{-1}L$ and
$A_L=\Lambda\otimes S_L$ in the original matrix coordinates.
$S_L$ is the source's ideal quotient since nonzero fractional
ideals are invertible. It is self-dual because its multiplying
ideal has product one with its conjugate. The tensor lattice is
also self-dual for the matrix trace pairing: a local basis of one
factor and its dual reduce this assertion to the duality in the
other. For two levels the common ideal factor cancels from the
local lengths, and $\Lambda$ is locally free of rank $n$.
Taking the sum of lengths at each prime and multiplying the local
indices gives
\[
          [A_L:A_L\cap A_{L'}]=[L:L\cap L']^n.
\]

Put $u_L$ equal to the normalized projection of the original
vacuum onto the $A_L$-vacuum line. For two self-dual tensor
lattices $A,B$, the inner product of their unit vacua has modulus
$[B:A\cap B]^{-1/2}$ with the prescribed common lattice
semicharacter. Indeed average the $B$-translations of the
$A$-vacuum with their inverse characters over $B/(A\cap B)$.
The two characters agree on the intersection. Distinct terms of
this average are orthogonal, as translation by an element of $A$
distinguishes two different cosets by the perfect pairing.
The average is the orthogonal projection onto the $B$-vacuum
line: it is invariant with the prescribed character, and its
inner product with a vector on that line equals the unprojected
inner product. Its squared norm is the reciprocal of the index.
This also proves that the projected vacua being normalized here
are nonzero. The two directions of the relative index coincide by
F74N1's dual-index argument.

In the finite sum defining $\Phi$, the first $r$ rows encode an
element of $J$ in the chosen residual frame and the remaining rows
are zero modulo $\mathfrak m$. Choose one matrix $y_k$ of representatives
for each $k\in J$, simultaneously for $L$ and $L'$. To justify
this, use the lattice with target factor
$2\bar{\mathfrak m}(L\cap L')$. At every prime of $\mathfrak m$
both marked lattices equal the reference lattice, the source
lattice equals the standard one, and $2\bar{\mathfrak m}$ is a
unit locally. Reduction of this full lattice onto the finite
quotient is surjective by the Chinese remainder theorem. The
specified residual frame is a basis there: it spans, and a
surjection between free modules of the same rank over each local
factor is invertible. These residues are therefore distinct.
In the defining sum, a change within a residue class lies in
the denominator lattice with target factor
$2\bar{\mathfrak m}\mathfrak mL$ (or $L'$). Its even lattice
translation fixes the vacuum: write the change as $2z$ with
$z\in A_L$, so its semicharacter is one. A numerator representative
has the form $2y'$; the pairing of $y'$ with $2z$ is integral
by tensor-lattice integrality. Thus the pairing of $2y'$ with $2z$
is even and the Heisenberg product phase is one. Equivalently one may
use the native common representative formula
$\Phi_L=\sum_{k\in J}\lambda(k)\,T_{y_k}u_L$, with exactly the
same $y_k$ for $L'$.

A nonzero cross inner product at $y_k,y_{k'}$ forces
$B(\eta,y_k-y_{k'})\in\mathbb Z$ (the sum of the row pairings)
for every $\eta\in A_L\cap A_{L'}$,
by applying the common translation and cancelling its character.
Thus $y_k-y_{k'}\in(A_L\cap A_{L'})^\vee=A_L+A_{L'}$.
The latter identity follows by dualizing the sum of the two dual
lattices. At the primes of $\mathfrak m$ this sum consists rowwise
of $\mathfrak m$ times the reference lattice. Hence $k=k'$.
For $k=k'$ unitarity of translation and $|\lambda(k)|=1$ leave
$\langle u_L,u_{L'}\rangle$. There are precisely $|J|$ such
terms, proving the two assertions, including the diagonal.
% END F74N2

% BEGIN F74N3
\paragraph{F74N3. The strict bound is uniform in the lattice and in the marking.}
Suppose $n\ge 4d+4$ in F74N2. The original norm field of
$\mathrm{IsSystem}$ now has diagonal $|J|>0$, majorant
$b(L,L')=[L:L\cap L']^{-n/2}$ and off-diagonal sum at most $0.65$.
Indeed for each finite set of lattices other than $L$, group by the
integer $m\ge2$ of F74N1. The sum of the majorant is bounded by
\[
   \sum_{m=2}^{\infty} m^{2d-n/2}
       \le \sum_{m=2}^{\infty}m^{-2}
       =\pi^2/6-1<0.65 .
\]
All terms are nonnegative, so this proves summability and the same
bound for the entire set. The argument takes the supremum over
finite sums and needs no ordering of the period lattices. The
constant is independent of the class representative $\Lambda$
and of $L$. In particular the diagonal term dominates the entire
off-diagonal sum even at the endpoint $n=4d+4$.

For the finite positive class weights of ICL12 put
$a=|J|\sum_{[\Lambda]}\mathrm{wt}(M,\Lambda)$.
The weighted matrix has diagonal $a$ and the off-diagonal row sum
is at most $0.65a$. On the source space with its prescribed
covariance, once the level action preserves that space, its inverse
is consequently given by the Neumann series with norm at most
$1/(0.35a)$. This supplies the quantitative norm hypothesis
previously assumed in ICL12 for these test vectors. The covariance,
support, cuspidality and doubled-family admissibility fields of
the full level system have independent content.
% END F74N3

\section{Reproduction by the original real Gaussian kernel}

Use the frame domain and the Lebesgue measure in ICL14. Write
$p=X.p_0$, $q=r-p$, $J=\operatorname{diag}(1_p,-1_q)$, and
$S_\delta$ for the open slab on which the Gram matrix at each CM
embedding differs entrywise from $J$ by less than $\delta$. Fix the
original finite family $P_b$ of harmonic row-highest polynomials of
the indicated weights, with at least one nonzero member. Write $K$ for
its actual doubly differentiated Gaussian kernel and $\lambda$ for
its Borel weight. We prove (ICL.7) on the entire domain. The
archimedean calculation needs even $n\ge 2r$, so includes the range
of ICL14.

% BEGIN K74A1
\paragraph{K74A1. The integrable determinant on a split Gram slab.}
At one place use ordinary complex matrix coordinates
$Y=\left(\begin{smallmatrix}A&B\\C&D\end{smallmatrix}\right)$
with $A$ of size $p\times p$. Entrywise conjugation interchanges the
source convention $Y^tJ\bar Y$ with $Y^*JY$ without altering volumes
or determinant norms. We claim that for sufficiently small $\delta$
the function $|\det A|^{-n}$ is integrable over the full slab
$Y^*JY\simeq J$ (here $\simeq$ means an entrywise error $<\delta$).

Put $U=\binom AC$, $G=U^*JU$. By shrinking $\delta$,
$G$ is uniformly near $I$ in operator norm, say
$\frac34 I<G<\frac54 I$. Conditional on $U$, each remaining column
$v$ satisfies bounds on $U^*Jv$ and $-v^*Jv$ by constants depending
only on $r$. Consequently $v$ lies in an ellipsoid
$v^*Hv\le R$ with $H=2JU U^*J-J$ and $R$ uniform. Indeed $H$ is
positive on the $J$-negative orthogonal complement of the columns
of $U$, where it equals $-J$. On the positive column space its
form is $2G^2-G$, which is positive. The spaces are orthogonal also
for $H$. The matrix determinant lemma gives
$|\det H|=|\det(I-2G)|$, uniformly bounded away from zero.
The real volume of this complex ellipsoid is its ball volume
divided by $\det H$. Dropping all cross-conditions among the last
$q$ columns bounds their joint volume by a uniform constant.

For the remaining $A,C$, positivity of $G=A^*A-C^*C$ makes $A$
invertible. Substitute $C=ZA$. The Jacobian is $|\det A|^{2q}$.
Set $P=I-Z^*Z$, which is positive since $G=A^*PA>0$. Thus $Z$
ranges in a bounded matrix ball. For fixed $Z$ set
$A'=P^{1/2}A$, so $G=A'^*A'$, and $A'$ ranges in a bounded set
with determinant bounded away from zero. The factors involving $P$
in the volume times $|\det A|^{-n}$ now multiply to
\[
                 (\det P)^{(n-2q)/2-p}.
\]
This is at most $1$ since $0<P\le I$ and $n\ge2(p+q)$. Integration
over the two bounded sets proves the claim, including the endpoint
$n=2r$. Empty matrix blocks give the same formulas with empty
determinant $1$. Left multiplication by an invertible fixed complex
matrix and permutation of coordinates change Lebesgue measure by a
constant positive factor. The claim therefore applies in any of the
original frame coordinates.
% END K74A1

% BEGIN K74A2
\paragraph{K74A2. Uniqueness of the lowest weight line.}
Fix an orthonormal $x_*$. For the product of unitary diagonal tori
let $a_t=x_*t x_*^{-1}$, $\chi(t)=\prod_{\tau,j}t_{\tau j}^{
\lambda_{\tau j}}$, with normalized Haar measure $dt$.
Two holomorphic functions of the given right weight and this left
torus character satisfy $G_1(y)G_2(x_*)=G_2(y)G_1(x_*)$. To prove
it, subtract the two sides and call the difference $m$. On the big
cell write $y=x_*Lb$, with $L$ lower triangular, $b$ upper triangular,
and use right covariance. Let
$D(s)=\operatorname{diag}(1,s,\ldots,s^{r-1})$ at each embedding.
Then $D(s)LD(s)^{-1}$ has entries $s^{i-j}L_{ij}$ for $i\ge j$,
and extends to the diagonal of $L$ at zero. The resulting frame
stays in the domain for $|s|\le1$. In the signature test on each
positive partial column space, the rows beyond the signature break
are multiplied by powers of $|s|$ no larger than those on the
positive rows. Hence their negative contribution cannot destroy
positive definiteness. The matrix stays invertible for $s\ne0$. Together
with positivity of the first $p$ columns this gives the original
domain condition. Right diagonal rescaling preserves it, and at zero
the lower-triangular factor becomes an invertible diagonal. On $|s|=1$ the left and right diagonal
characters cancel, so $m(x_*D(s)LD(s)^{-1})=m(x_*L)$.
The function is holomorphic on the disk and continuous on its
boundary. It is constant, and its center value is zero by right
covariance and $m(x_*)=0$. This proves vanishing on the big cell.
Frames with nonzero leading minors are dense in the open domain
(perturb by a small scalar multiple of $x_*$), so $m=0$ throughout.
For a function with nonzero value at the orthonormal frame, the
same calculation gives its normalized value $1$ on the lower
unipotent cell.
% END K74A2

% BEGIN K74A3
\paragraph{K74A3. The actual Gaussian, its weight and positive diagonal.}
At frames $x,y$ the kernel is obtained by applying the constant
coefficient operators $P_b$ in $z$ and $\overline{P_b}$ in $w$
to the doubled Gaussian at $z=w=0$, and summing over $b$. In the
doubled-frame expression, split the quadratic exponent into its
$z,z$, $w,w$ and $z,w$ blocks. Within each of the first two blocks
the only quadratic factors are sums of the form
$\sum_a Z^+_{ia}Z^-_{ja}$ (and their $w$ counterparts), with
coefficients independent of the column $a$. This follows by writing
the frame in its positive and negative row blocks: the same-sign
diagonal blocks of the pure quadratic correction vanish, and the
off-diagonal blocks pair the two row types. In the derivative pairing,
multiplication by this quadratic is adjoint to
$\sum_a\partial^+_{ia}\partial^-_{ja}$. Harmonicity annihilates
every positive power of each pure correction, so these exponentials
can be removed in computing the derivative. The mixed exponential
and its determinant prefactor are holomorphic in the first frame.
Indeed their entries are obtained by inverting the transverse
mixed frame matrices, whose determinants do not vanish on the
domain. Right triangular multiplication substitutes its positive
and negative blocks in the two sets of letters. The highest-weight
equations for $P_b$ and the determinant factor give precisely
$K(xb,y)=\operatorname{bchar}_\lambda(b)K(x,y)$.
Interchanging $x,y$ conjugates the mixed Gaussian and exchanges the
two derivative operators. An isometry of the Hermitian space on
the left of both frames preserves their mixed Gram expressions,
and its determinant has modulus one. Hence
$K(y,x)=\overline{K(x,y)}$ and $K(ux,uy)=K(x,y)$.

For an orthonormal frame $x_*$ the mixed exponential is
$\exp(\sum_l t_l z_l w_l)$ in the dual letter coordinates, with
$t_l=-\pi\,\mathrm{Im}(\tau_l\kappa_0)>0$. The prefactor is
$D_* =\prod_\tau |\det x_{*,\tau}|^n>0$. Expanding the
exponential and differentiating gives
\[
      K(x_*,x_*)=
        D_*\sum_{b,\alpha}\alpha!\,
           |(P_b)_\alpha|^2\prod_l t_l^{\alpha_l}>0 .
\]
This proves positivity without assuming that a chosen theta section
has nonzero value at any frame.

For later use we describe why $K(x,\cdot)$ is both bounded and
integrable on all sufficiently thin $S_\delta$, with $x$ fixed.
First take $x=x_*$ and an orthonormal $y$. In coordinates
$Y=x_*^{-1}y$ the holomorphic function $K(y,x_*)$ has right Borel
weight $\lambda$ and the corresponding left torus weight, by
invariance and symmetry above. Normalize its value at the identity.
On the big cell its value is the product of the leading minors
$D_j(Y)$ with exponents $\lambda_j-\lambda_{j+1}$, including the
last determinant with exponent $\lambda_r$. Indeed on the lower
unipotent cell its normalized value is $1$ by the contraction of
K74A2, and the upper triangular diagonal in Gaussian elimination
is $D_j/D_{j-1}$. Equalities extend to bounds at all orthonormal
frames by continuity. For $Y^*JY=J$, $|D_p|\ge1$,
$|D_r|=1$ and $|D_j|\le |D_p|$ for all $j$. To see the last
inequality, exchange the last $q$ output coordinates with the
last $q$ input coordinates in the isometry $Y$. Conservation of
the indefinite form gives a unitary map on the resulting positive
definite coordinates. Its appropriate principal minors are
$D_j/D_p$ up to sign or conjugation, by block elimination, so their
norms are at most one. Equivalently for $j\le p$, the upper $p$
block has all singular values at least one and its $j$-minors are
bounded by $|D_p|$. For $j>p$ apply the complementary-minor formula
and the same reasoning to the lower $q$ block of $Y^{-1}$.

The weights are decreasing within either signature block, and across
the break $\lambda_i+n\le\lambda_j$ for $i\le p<j$.
Bounding all the nonnegative minor exponents and retaining the
one at the break therefore gives $|K(y,x_*)|\le
C_*\prod_\tau |D_p(Y_\tau)|^{-n}$. In the definite cases the
orthonormal set is compact and this bound also holds. On the thin
slab signed Gram--Schmidt writes $y=y' b$ with $y'$ orthonormal
and the diagonal entries of $b$ and their inverses uniformly
bounded: the Gram minors are close to those of $J$ and their
successive ratios give the squared diagonal moduli. Right Borel
covariance and $D_p(Y'b)=D_p(Y')\prod_{j\le p} b_{jj}$ extend
both the bound and this integrable majorant to the slab. K74A1
applies at each of the finitely many places and Fubini gives
integrability on the product. Every fixed $x$ is an orthonormal
frame times a right Borel factor, and left unitary transport preserves
the slab and Lebesgue measure. Both assertions follow for this $x$.
% END K74A3

% BEGIN K74A4
\paragraph{K74A4. The reproducing scalar on the original slab.}
Choose $\delta_1>0$ small enough for K74A3. For every
$0<\delta\le\delta_1$, choose an orthonormal $x_*$. Put
\[
   \rho_\delta=
   \frac{\int_{S_\delta}|K(x_*,y)|^2\,dy}{K(x_*,x_*)}.
\]
The denominator is positive by K74A3. The numerator is finite by
the simultaneous integrability and boundedness there. It is
positive because the continuous integrand is positive at $x_*$
and the slab is an open neighborhood of that point. These arguments
include a zero-dimensional frame space, with its ordinary Haar
measure.

For a bounded
holomorphic form $F$ of the given right weight set
\[
                  G(y)=\int \overline{\chi(t)}F(a_ty)\,dt.
\]
The map $G$ is holomorphic: on compact subneighborhoods the
translated arguments range over a compact subset of the domain, so
local uniform bounds justify integration of the holomorphic
functions. It has the same right weight, left torus character
$\chi$, and $G(x_*)=F(x_*)$ by right covariance. For fixed
$k(y)=K(x_*,y)$ one has $k(a_ty)=\chi(t)^{-1}k(y)$.
The $a_t$ preserve both the slab and Lebesgue measure (their complex
determinants are units of modulus one). A form bounded in its
weighted norm is bounded on the slab by the Borel diagonal bounds
of K74A3. Thus Fubini and a change of variables give
$\int_{S_\delta} kG=\int_{S_\delta} kF$.

The function $y\mapsto K(y,x_*)$ has the two transformation laws of
K74A2 by K74A3. Apply uniqueness there with the torus mean $G$:
\[
   G(y) K(x_*,x_*)=K(y,x_*)F(x_*).
\]
Multiply by $K(x_*,y)$ and integrate. Symmetry of $K$ and the
torus-mean identity give the reproducing formula at $x_*$. For
any orthonormal $x'=u x_*$ apply this formula to $F(u\,\cdot)$.
Left multiplication by $u$ preserves the kernel, slab and measure.
For an arbitrary domain frame $x=x'b$ apply the right Borel law to
both sides, with the same factor. Thus for every bounded holomorphic
form with the specified Borel weight and every domain frame,
\[
             \int_{S_\delta} K(x,y)F(y)\,dy
                          =\rho_\delta F(x).
\]
The same strictly positive scalar works for all $F,x$ at this
width. The slab, measure, and kernel in this formula are those of
(ICL.7). The separate arithmetic identity (ICL.6) involves a sum
over the rational unitary group, a class sum, and the original
vacuum multipliers. That identity and the geometric level-system
construction remain necessary to apply ICL14.
% END K74A4

\section{The arithmetic coefficients of the doubled kernel}

% BEGIN O74E1
\paragraph{O74E1. The graph indexing set.}
Let $X$ be the source Hermitian datum over the CM field $E$, with
nonsingular Hermitian form $h$ on $E^r$, Gram matrix $G$, and
polarization parameter $\kappa_0\ne0$. The indices of its doubled
Eisenstein series are the bottom row spaces $U$ of the rational
split unitary matrices $\sigma$ of size $2r$. Put
\[
 T=\begin{pmatrix}1&(\kappa_0/4)G\\-1&(\kappa_0/4)G\end{pmatrix},
 \qquad i_\gamma=T^{-1}\operatorname{diag}
       (1,(\gamma^{-1})^{t})T .
\]
For every $U$, exactly one of the following holds. Either $U$ is the
bottom row space of $i_\gamma$ for a unique unitary matrix $\gamma$ with entries in $E$,
or $U$ contains $(\xi,0)T$ for some nonzero $\xi$ with $h(\xi,\xi)=0$.

Indeed, $T$ is invertible. Multiplication in blocks gives the pairing
on the inverse image $I=UT^{-1}$ as $(\kappa_0/2)(h\oplus(-h))$.
The bottom rows of $\sigma$ are independent and mutually isotropic,
so $I$ is isotropic of dimension $r$. If projection to the second
summand has a nonzero kernel, any nonzero $(\xi,0)$ in it works.
Otherwise this projection is an isomorphism and $I=\{(Dv,v)\}$.
Isotropy gives $h(Dv,Dw)=h(v,w)$ for all $v,w$. Nondegeneracy makes
$D$ injective and hence a unitary automorphism. To identify the
index with the source matrix, multiply a row $(0,c)$ by $i_\gamma$.
The result is $(u,\gamma^{-1}u)T$, where
$u=c((\kappa_0/2)G)^{-1}$ ranges through all rows and
$\gamma^{-1}$ acts on the coordinate vector as a column.
Choosing $\gamma=D$ gives the claimed space and recovers $\gamma$
uniquely. Its projection to the second summand is injective.
% END O74E1

% BEGIN O74E2
\paragraph{O74E2. A common nonzero scalar at the origin.}
Fix the symplectic basis of $X.\Lambda$, the labeling
$\{1,\ldots,g\}\simeq\Psi\times\{1,\ldots,r\}$ and even $n$.
Write $D(\phi,\psi)$ for the native doubled test function.
Use its theta product formula, sesquilinearity and Heisenberg
intertwining as inputs here. There is a single $a\in\mathbb C^\times$
for all rational Schwartz--Bruhat tests such that
\[
 B(\phi,\psi):=D(\phi,\psi)(0)=a\langle\phi,\psi\rangle,
 \qquad \langle\phi,\psi\rangle=\int\phi\overline\psi .
\]
In rational symplectic coordinates put $d=ng$. The simultaneous action
of a rational Heisenberg displacement on both arguments becomes a pure
modulation on $D$: the translation block is $y-y=0$. Its value at zero
is one as a multiplier. Every translation and modulation on
$\mathbb Q^d$ occurs because the chosen basis identifies the rational
phase space with $\mathbb Q^{2d}$. Thus $B$ is invariant under both
simultaneous actions.

At a common positive level $R$, expand the tests in the indicators
$e_u$ of the cosets $u+R\mathbb Z^d$ inside $R^{-1}\mathbb Z^d$.
On $e_u$, modulation by the $j$th coordinate divided by $R$ acts by
$\exp(2\pi i u_j/R)$. For distinct cosets some $j$ distinguishes the
two characters. Invariance and sesquilinearity give $B(e_u,e_v)=0$
for $u\ne v$. Translation makes all diagonal values equal. The
source Haar normalization gives mass $R^{-d}$ to each such coset,
so the proportionality scalar at this level is $R^d B(e_0,e_0)$.
The test $u_0=\mathbf1_{\mathbb Z^d}$ has norm squared one and belongs
to every level, hence this scalar always equals $B(u_0,u_0)$.
Any two tests have a common level by taking a multiple of their levels.

To prove $a\ne0$, first $D(u_0,u_0)\ne0$. Otherwise evaluate the
theta product identity on equal frames with the second elliptic
argument the conjugate of the first. It gives the absolute square
of $\Theta_{u_0}$ identically zero. At any regular frame the period
change is invertible, the Gaussian prefactor is nonzero, and the
Fourier coefficient at frequency zero of the ordinary theta series
is $u_0(0)=1$. This is impossible. If $a=0$, apply $B$ to the
first test translated by the native displacement $-t$ and to $u_0$
in the second slot. Doubling sends this displacement to the Hermitian
Heisenberg operator with translation $-t$. Its value at zero is
$D(u_0,u_0)(t)$ times a nonzero phase for every
$t\in\operatorname{Mat}_{n,r}(E)$. The assumed vanishing would give
$D(u_0,u_0)=0$, a contradiction.
% END O74E2

% BEGIN O74E3
\paragraph{O74E3. The singular terms and the surviving coefficient.}
Suppose $4r<n$, $M>0$ and an actual level system
$(\Delta,\mathcal L,\Phi,\mu)$ is given in the sense of
\texttt{IsSystem}. Fix $L,L'\in\mathcal L$, an admissible lattice
$\Lambda\in V(E,n,M)$, and frames $x,y$ in the domain of $X$.
Put $\phi=\Phi(L,\Lambda)$, $\psi=\Phi(L',\Lambda)$.
For independent elliptic variables $(z,w)$ let $p$ denote the
doubled variables $(\mathrm{dblZ},\mathrm{dblz},\mathrm{dblz}')$
at $(x,y;z,w)$ and $k=k_0(x,y;z,w)$ the native Gaussian prefactor.
Use the source Weil and theta transformation laws and the graph
coordinate change. The contribution of $U$ to $k\,\mathrm{Eis}(D(\phi,\psi);p)$
is zero unless $U$ is the graph of an element $\gamma\in\Delta$.
For such a graph its entire contribution is
\[
 a\,\langle\Phi(L,\Lambda),\Phi(\gamma L',\Lambda)\rangle
       \,\overline{\mu(\gamma,L')}\; k_0(x,\gamma y;z,w).
\]

For the singular case of O74E1 use the system's cusp condition on
$\xi$ and $\phi$. It supplies a finite combination
$\sum_v b_v H_G(v\xi^t)\phi=0$ with $\sum_v b_v\ne0$.
Write $U$ as the bottom row space of $\sigma$. Every row of
$(v\xi^t,0)T$ belongs to $U$, so the first block of
$(v\xi^t,0)T\sigma^{-1}$ is zero. The Weil covariance and the
doubling law therefore replace each displacement by pure modulation
after applying $\omega(\sigma)$ to the doubled function. Evaluation
at zero is unaffected. Linearity in the first test gives
$(\sum_v b_v)(\omega(\sigma)D(\phi,\psi))(0)=0$.
This kills the numerator of the Eisenstein summand.

For a graph one may use $i_\gamma$ as representative. Indeed if
$\sigma,\sigma'$ have the same bottom space, $P=\sigma'\sigma^{-1}$
has lower left block zero. Write its blocks as $(A,B;0,D_0)$, so
$D_0=(A^*)^{-1}$. Its factor of automorphy is the constant
$d_0=\prod_{\tau\in\Psi}\tau(\det D_0)^n\ne0$. At the origin its
Weil action also multiplies by $d_0$. To check the latter formula,
the defining theta transformation for $P$ changes a summation matrix
$t$ into $tA$ and multiplies the test by $d_0$ and the quadratic
character with matrix $-BA^*$. In fact the new period is
$(AZ+B)A^*$, whose extra term $BA^*$ cancels that character, while
the linear elliptic terms are carried by the same change of index.
At $t=0$ the character equals one. The Weil composition law and the
cocycle of $\mathrm{hFac}$ thus multiply numerator and denominator
by the same $d_0$ under $\sigma\mapsto P\sigma$.

Write $p_\gamma,k_\gamma$
for the same variables and scalar with $y$ replaced by $\gamma y$,
and $j_\gamma=\mathrm{hFac}(i_\gamma;p)$. The native graph coordinate
law gives $i_\gamma p=p_\gamma$ and $k_\gamma j_\gamma=s_\gamma k$
for a nonzero constant $s_\gamma$. The defining theta product and the
ordinary unitary covariance $\Theta_{w_G(\gamma)\psi}(\gamma y)=\Theta_\psi(y)$
give
\[
 \Theta_H(\omega(i_\gamma)D(\phi,\psi);p_\gamma)
  =s_\gamma\Theta_H(D(\phi,w_G(\gamma)\psi);p_\gamma).
\]
Indeed the left side is $j_\gamma$ times the original doubled theta,
and multiplication by $k_\gamma$ uses precisely the displayed scalar
law. Varying $(z,w)$ gives all doubled elliptic variables. At fixed
period, the distinct rational matrix frequencies in $\Theta_H$ have
nonzero Gaussian factors and separate the coefficients by Fourier
uniqueness on a common period torus. Thus the tests themselves agree
up to $s_\gamma$. The graph summand is
$k(\omega(i_\gamma)D)(0)/j_\gamma
 =k_\gamma D(\phi,w_G(\gamma)\psi)(0)$.
By O74E2 the latter test value is $a\langle\phi,w_G(\gamma)\psi\rangle$.
The system's support law makes this zero for $\gamma\notin\Delta$.
Inside $\Delta$ the second test is $\mu(\gamma,L')\Phi(\gamma L',\Lambda)$.
Conjugate linearity in that slot gives the stated conjugation.
% END O74E3

% BEGIN O74E4
\paragraph{O74E4. The weighted sum and its elliptic derivatives.}
As convergence input it suffices to take absolute summability of the
native Eisenstein series at \emph{zero} elliptic variables in the
stable range. Indeed fix its period $Z$. At one place let $C,D_0$
be the bottom blocks of a representative, $A=CZ+D_0$, and
$B=A^{-1}C$. The unitary row relation $CD_0^*=D_0C^*$ gives
\[
 -\operatorname{Im}B=B(\operatorname{Im}Z)B^*.
\]
This follows by multiplying $B-B^*$ on the left by $A$ and on the
right by $A^*$. If $\operatorname{Im} Z\ge\delta I$, it gives
$\delta\|B\|^2\le\|\operatorname{Im}B\|\le\|B\|$ in operator norm.
A common $\delta>0$ works at the finitely many places, so for fixed
doubled elliptic variables the inverse Gaussian factor in each term
is uniformly bounded. Explicitly at one place the bound is
$\exp(2\pi\delta^{-1}\sum_k\|\zeta'_k\|\|\zeta_k\|)$.
Both the determinant and the finite test value are unchanged when
these variables are set to zero. The bound therefore transfers
absolute summability to every $(z,w)$. For $M>0$ use the finite class set of
the level $M$ and its chosen representatives in the definition of
$\operatorname{classSum}_M$. A representative has weight equal to the
reciprocal of the order of its stabilizer in the source group $D_H(M)$.
Set
\[
 c_\gamma=a\,\overline{\mu(\gamma,L')}
 \operatorname{classSum}_M\bigl(\Lambda\mapsto
  \langle\Phi(L,\Lambda),\Phi(\gamma L',\Lambda)\rangle\bigr).
\]
By O74E1 and O74E3, extending the graph terms by zero gives the
original summands for every doubled isotropic space. Absolute
summability allows restriction to this injectively indexed family.
The finite class sum commutes with it. Consequently
\[
 S(z,w):= k\,\operatorname{classSum}_M\bigl(\Lambda\mapsto
   \mathrm{Eis}(D(\Phi(L,\Lambda),\Phi(L',\Lambda));p)\bigr)
    =\sum_{\gamma\in\Delta} c_\gamma k_0(x,\gamma y;z,w).
\]
In particular the stabilizer weights belong to the definite lattice
classes. Each group element labels one graph, including when distinct
elements act identically on $L'$.

The series may be differentiated termwise to every finite order in
both elliptic variables. To see this, put $v=(z,w)$. The source
formula for $k_0$ has a constant determinant prefactor and the
exponential of a homogeneous complex quadratic $Q_\gamma(v)$, since
both doubled elliptic variables are linear in $v$. Write a term as
$A_\gamma\exp Q_\gamma(v)$. In the real dual of the finite-dimensional
space of complex homogeneous quadratics choose a basis from the
functionals $e_u:Q\mapsto\Re Q(u)$. They span the dual, since their
common kernel is zero (use $e^{\pi i/4}u$ to detect the imaginary
part). Their negatives are $e_{iu}$. For compact $K$, the coordinates
of $e_v$ in that basis have uniformly bounded sum of absolute values
as $v$ ranges over $K$. Scaling all basis vectors by a common constant
therefore puts $e(K)$ inside the convex hull of finitely many $e_{u_j}$:
a positive multiple of $e_u$ equals $e_{\sqrt{t}u}$ for multiplier $t$.
In zero dimension use the single probe $u_1=0$. It follows that
\[
 \sum_\gamma\sup_{v\in K}|A_\gamma e^{Q_\gamma(v)}|
 \le \sum_j\sum_\gamma |A_\gamma e^{Q_\gamma(u_j)}|<\infty .
\]
Cauchy's formula on a larger polydisc gives the same convergence for
each prescribed derivative on a smaller one. This proves the claim.

In particular fix $p_0+q_0=r$, row weights $c,c'$, and a finite
family $(P_j)$ of the corresponding row-highest harmonic polynomials
satisfying the source unitary column-invariance condition
\texttt{InvFam}. Let $\mathcal D_P$ be the sum over $j$ of
$P_j(\partial_z)\overline{P_j}(\partial_w)$ at zero in the original
letter directions. Since $K_P=\mathcal D_P k_0$, the differentiated
Eisenstein side is $\sum_\gamma c_\gamma K_P(x,\gamma y)$. To obtain
(ICL.6) one still needs to replace the Eisenstein side by the theta
side. Precisely, put
\[
 H(Z,\zeta,\zeta')=\operatorname{classSum}_M
   \bigl(\Lambda\mapsto(\Theta_H-\mathrm{Eis})
       (D(\Phi(L,\Lambda),\Phi(L',\Lambda));Z,\zeta,\zeta')\bigr).
\]
For the pair $L,L'$ the system's admission clause supplies $R$ for
this doubled family. O74E5 below proves the native heat equations for $H$. The
remaining two inputs here are the definite Siegel--Weil identity
$H(Z,0,0)=0$ for every allowed $Z$ and the invariant-polynomial
heat identity that these together imply $\mathcal D_P(kH(p))=0$ when $r\le n$. Under these
inputs the theta product formula identifies $S+kH(p)$ with the finite
weighted sum of
$\Theta_{\Phi(L,\Lambda)}(x;z)\overline{\Theta_{\Phi(L',\Lambda)}(y;\bar w)}$.
The real letter directions and conjugated polynomial in the second
slot give exactly the weighted sum of the differentiated theta
products in (ICL.6). The scalar $a$ was fixed in O74E2 before choosing
any level system, polynomials, lattices, or frames. Existence of the
level system, the stated Weil, theta and convergence inputs, and the
two remaining assertions about $H$ retain their separate hypotheses.
% END O74E4

% BEGIN O74E5
\paragraph{O74E5. The heat equations for the actual weighted difference.}
The function $H$ of O74E4 is holomorphic on the open product of the
Hermitian half spaces and the two elliptic matrix spaces, and satisfies,
at each embedding $\tau$ and for every $a,b$,
\[
 \lambda\frac{\partial H}{\partial Z_{ab}}
  =\sum_{k=1}^n\frac{\partial^2 H}
       {\partial\zeta_{ka}\,\partial\zeta'_{kb}},\qquad \lambda=2\pi i.
\]
All derivatives are complex coordinate derivatives at $\tau$ with
the other embeddings fixed. The two elliptic matrices are independent.
We use the zero-argument stable-range convergence input of O74E4 at
every allowed period. This suffices also for locally uniform convergence
in the period.

For the theta side write $u_k$ for the row of the summation matrix at
$\tau$. Its exponential there is
$\exp\{\lambda\sum_k(u_k Z\bar u_k^t+u_k\zeta_k+
 \zeta'_k\bar u_k^t)\}$, with $\zeta_k$ a column and $\zeta'_k$ a row.
Each side of the differential equation multiplies the summand by
$\lambda^2\sum_k u_{ka}\overline{u_{kb}}$.
To justify the derivatives, on a small closed polydisc each
$\operatorname{Im} Z_\tau$ has a common positive lower bound and both
elliptic matrices are bounded. The modulus of a term is then bounded
by a constant times $\exp(-c\|u\|^2+C\|u\|)$ with $c>0$, taking
all embeddings of the CM type together. The fixed finite test is bounded
and supported in a fixed denominator lattice in Minkowski space by
its Schwartz--Bruhat level. A Gaussian with this linear error is
summable over that lattice, giving local normal convergence.

For an Eisenstein summand keep the notation $A=CZ+D_0$ and
$B=A^{-1}C$ at the given place. Its variable factor is
\[
 F=(\det A)^{-n}\exp\left(-\lambda\sum_k\zeta'_k B\zeta_k\right).
\]
For an increment $V$ in $Z$, the differential of $B$ is $-BVB$,
and the logarithmic differential of $\det A$ is $\operatorname{tr}(BV)$.
Take $V$ to have its only nonzero entry $1$ in position $(a,b)$.
Dividing $\lambda\partial_{Z_{ab}}F$ by the nonzero factor $F$ gives
\[
 -n\lambda B_{ba}+\lambda^2\sum_k(\zeta'_k B)_a(B\zeta_k)_b .
\]
The mixed derivative for a single $k$, divided by $F$, contributes
$-\lambda B_{ba}+\lambda^2(\zeta'_kB)_a(B\zeta_k)_b$.
Summing gives the same expression. The test value at zero and all
other archimedean factors are constant for these derivatives.

Fix a central period $Z_*$ and put $B_*=(CZ_*+D_0)^{-1}C$ at each
place. If $\operatorname{Im} Z_*\ge\delta I$, O74E4 gives
$\|B_*\|\le\delta^{-1}$ uniformly in the representative. For
$Z=Z_*+V$ with $\|V\|<\delta/2$, factor
\[
 CZ+D_0=(CZ_*+D_0)(I+B_*V).
\]
All singular values of the second factor are at least $1/2$.
The inverse determinant to power $n$ grows by at most $2^{rn}$
at each place. Also $\operatorname{Im} Z\ge(\delta/2)I$; the
elliptic exponential has a uniform bound on every bounded set of
elliptic arguments by O74E4. The absolute terms at $(Z_*,0,0)$
therefore supply a single summable majorant, up to a constant,
throughout a neighborhood of any prescribed point. Cauchy's formula
now allows the holomorphic sums and their derivatives to be taken
termwise. The finite weighted difference over the definite lattice
classes preserves the equation and holomorphy. This gives precisely
\texttt{IsHeat} for the original $H$.
% END O74E5

\section{The base rows and the theta square on a single linked bundle}

% BEGIN B74R1
\paragraph{B74R1. A single witness for the dependent data.}
Fix a parameter $p$, an original weighted point $q$ and a bundle $b$ at
that point, and suppose $w:\mathrm{BWit}(q,b)$. Put $u=w.x$,
$R=u.R$, $F=u.F$, $Y=u.Y$, $Z=u.Z$, $X=R.\mathrm{toCx3}(F)$.
The proof field of $w.y$ is a witness of $\mathrm{Holds}(R,F)$.
The equality in that same record identifies two elements of
$\mathrm{RPk}(b.R,b.S)$:
\[
 (b.P,b.C.c.\mathrm{face},b.C.\mathrm{Sp},\mathrm{zsc}_b,b.B)
 =
 (P_u,F.\mathrm{Dc},F.\mathrm{Sp},F.\mathrm{Zsc},
            \mathrm{realTBOf}(w.y.y)).
\]
Here $P_u$ is the row-side family built from $u$, and $\mathrm{zsc}_b$
uses the original pair of integers, scale and row-side family of $b$.
Any proposition depending on these five entries can be proved on the
right and transported to the left by substitution. The later entries
depend on the earlier ones, so projection to the first coordinate
alone would lose information. In the right-hand description the
realization is $R$, its face is $F$, and its chosen subspaces and
base rows are $X.\mathrm{zs}$ and $X.\mathrm{data}(F.\mathrm{cr},
\mathrm{false})$. All four identifications follow by expanding the
construction. The weak characterization comes from $u.\mathrm{pins}$,
whose phase field gives $R.\mathrm{phs}(j)=1$ for every lift $j$.
If one starts with the selected argument $w.a.x$, use $w.ax$ first.
This orders the identifications without choosing new column
spaces, bases or a second tuple reading.
% END B74R1

% BEGIN B74R2
\paragraph{B74R2. Base tests and smooth extensions for the stored map.}
Use B74R1 and write $K=D.\mathrm{what}_2.\mathrm{LL}$ and
$\sigma=A_3.A.v_0.\mathrm{embedding}$. For every lift $j$ there
are finitely many $\psi_a\in X.W(j)^\vee$ with
\[
 \begin{split}
 &\bigcap_a \ker\psi_a=0,\\
 &\psi_a\bigl(X.\theta'_{jt}(X.jL_{jt}v)\bigr)\in\sigma(K)
               \quad(t\text{ arbitrary},\ v\in A.FL(j,t)),\\
 &\psi_a\bigl(X.w0C_{0j}(e_k)\bigr)\in\sigma(K)
               \quad(1\le k\le X.dW(0,j)).
 \end{split}
\]
In the last line $0$ denotes the base level and $e_k$ runs through
its stored coordinate vectors. Indeed, take
$w.y.\mathrm{intersectionNumber\_baseChange}.\mathrm{prim}$ and
expand its $\mathrm{liftTestB}$ field. The quantifier for $\psi$
precedes both rationality clauses. This proves separation on the
entire target, with the tests in the twisted model used in BTH8.
It retains the distinction between $\sigma$ and the distinguished
embedding used by the semilinear geometric reading.

The same $\mathrm{Holds}$ has an auxiliary field
$R.\mathrm{toRealDat}.\mathrm{ExAt}$. Its $\mathrm{thF}$ entry gives,
for every $j,t$, a map satisfying $\mathrm{PThF}$. Apply the first
implication of the weak choice rule in $u.\mathrm{pins}.\mathrm{char}$
to that witness. The \emph{stored} $U=R.\mathrm{thF}(j,t)$ then
takes every vector into the smooth cuspidal multiplicity space
$\mathrm{Wall}_j$ and restricts on pure tensors of lowest type to
the prescribed integral on $\operatorname{im}(R.\mathrm{extV}(j,t))$.
In particular evaluating $U$ on one basis at the base tuple gives the
finite extensions of BTH4. The analogous $\mathrm{thG}$ implication
identifies the stored projected map via $e_W\,\mathrm{thG}
=\mathrm{postOn}(eWf)\,U$. These deductions use the given
$\mathrm{Holds}$ only downstream of its construction.
% END B74R2

% BEGIN B74R3
\paragraph{B74R3. Constructing the eight base row fields.}
On the realization of B74R1 put $H=w.y.\mathrm{intersectionNumber\_baseChange}$.
Use $s=1$, $H.\mathrm{ok}$, $H.\mathrm{sep}$ and
$h=H.\mathrm{prim}$. We describe the base-level construction
$\mathrm{primRows3}$ with underlying data $X.\mathrm{data}(F.\mathrm{cr},0)$.
The weak model gives the rational forms and the row and lift isometries.
From $h.\mathrm{hecke}$ use the chosen row and lift operators together
with their intertwining equations. Use these very lift operators in
$h.\mathrm{liftOnto}$. The row and lift tests give rationality of the
good subspaces via rational orthogonal complements. Apply the native
unitary eigensystem decomposition with these subspaces, the rational
adjoints, lower-level type, unitary equivariance, the decomposition
into lift parts and the preceding surjectivity. Both outputs use one
eigensystem choice. They give the base types (using separation) and
the base proviso (using the given Hecke conditions).
The primitive positivity supplies the positive level, and the tuple
tie supplies the height identity $H0$.

For the frames, the row model gives the stated row frame, the lift
frame has invertible Gram matrix by the rational form, and the
integral lattice assertion is $h.\mathrm{mLat}$. The source's
base lift identity identifies the complexification of the rational
lift matrix with the theta map on the cut coordinates. Its rational
entries follow from the base lift isomorphism onto its image, the
principal frame and the tests of B74R2. The tests embed the target
in a coordinate space. The test matrix on the independent frame
columns has an invertible maximal minor over the embedded field.
Solving in those rows expresses each tested lift column over that
field. This is the native base rationality argument before forming
the row record. On the row side
the base row matrix is the identity, so its image of each type
subspace is precisely its prescribed rational part. On the lift side
the type membership test and the lift identity give the forward
inclusion. Lift-matrix surjectivity gives a preimage of each vector of
the target part, and the reverse membership test puts that preimage
in the specified source part. This gives both images in the field
called $\mathrm{l1a}$. The row and lift matrix bijections, scalar
positivities and the row and lift scalar transformation laws are the
entries of $\mathrm{l1c}$.

Here is the dimension argument for the last field. For each type $p$
let $a_p$ be the dimension over $K$ of its base lift coordinate
subspace, $b_p$ that of the corresponding rational target part, and
$c_p$ the complex dimension of the analytic part intersected with
the image of the base frame. The rational lift matrix embeds the
first subspace in the second. An independent family in rational
coordinates remains independent after scalar extension, and the
complex frame is injective. Hence $a_p\le b_p\le c_p$.
The complex intersections are independent by orthogonality, so
$\sum c_p$ is bounded by the dimension of the image of the frame.
The base lift is injective with precisely that image after extension.
Its source dimension is at most $\sum a_p$ since the type spaces
span the base coordinates. Thus equality holds separately for every
$p$, including a zero-dimensional or missing part. The row version
uses the base row image and its orthogonal parts by the same dimension
bounds. The primitive complex equality $h.\mathrm{dimParts}$ now
identifies the two rational part dimensions. This is $\mathrm{l1d}$.

Together positivity, proviso, types, $H0$, frames, the two image
equalities, the basis and scalar identities, and the dimension
equalities are precisely the eight additional fields of
$\mathrm{IndefRows}$. Its underlying data is the specified $X.\mathrm{data}$
by construction. Weak and primitive data also give stability and
the in-block assertion. Transport stability, block membership and
the row witness along the realization's equality for its subspaces,
then use its dependent equality for $S_0$. With $\langle 1,H.\mathrm{ok}\rangle$
and $H.\mathrm{sep}$ these are the five entries of
$\mathrm{RowsHold0}_3(b.P)$.
% END B74R3

% BEGIN B74R4
\paragraph{B74R4. The square and rank on every original linked selection.}
The same $Y$ contains $\mathrm{RecTied}$. Its lift field evaluates
the geometric theta map on a used tuple whenever a multiplicity
vector realizes the prescribed lowest-type integral and its good
projection belongs to the target. Here $Y.\mathrm{tup}(t_0)=\mu_0$
is used. The source and target Hilbert readings in $Y$, the coordinate
law $Z.\mathrm{jW\_read}$, and the fields
$w.y.r.\mathrm{mtRead},w.y.r.\mathrm{mjRead}$ all refer to these maps.
Furthermore $u.\mathrm{div}$ supplies the independence, spanning
and lattice identification of the target basis. These data belong to
the row-side family and realization identified in B74R1.
B74R2 supplies the smooth extensions and the simultaneous base
tests. BTH7 therefore evaluates the stored theta map on the full
level, BTH10 kills the orthogonal complement of the good source,
and BTH11 reads its restriction in the stored target coordinates.
Writing a coordinate vector as $E_j\Pi_jv+(v-E_j\Pi_jv)$ gives
BTH12's square
\[
 D.\theta P\bigl(\mu_0,j+1,D.jMt(\mu_0,j+1)(Y.eMt(j,t_0)v)\bigr)
       =Z.jW(j)(M_j\Pi_jv).
\]
Here $M_j$ is the base theta matrix, and $E_j,\Pi_j$ are the
inclusion and normalized projection for $b.P$. B74R3 supplies the
row witness to the base-rank argument SLR4--SLR6 on this very bundle.

The quantifiers give a stronger simultaneous conclusion on linked
selections. Fix any original class $\mathcal A$, raw selector $\sigma_b$
and row condition $E$ for which the source's
$\mathrm{IsSel}(Q,\mathrm{IsLinked}(p),E,\mathcal A,\sigma_b)$
holds. For each member $\mathcal F$, each $x>\mathcal F.x_1$,
and each index at $x$, its defining clause gives
$\mathrm{Built}$ for $b=\sigma_b(\mathcal F,x,i)$.
Eliminate the dependent equality in $\mathrm{Built}$, then the one
in $\mathrm{linkOf}$ defining $\mathrm{IsLinked}$. What remains
is precisely a weighted point $q$ and a witness $w$ for that
same $b$. Apply B74R1--B74R3 and the square above before transporting
the conjunction back. Thus all four BTH12 duties hold on this
selection at \emph{every} admissible index. The original eventual
rank statement follows with $x_0=\mathcal F.x_1+1$, independent of
the point and the lift. Taking $E$ identically true also proves the
source's universal selected-row implication ($\mathrm{Row.of\_sel}$).
It can be conjoined with the other row duties of DEX8.

The existence step in DEX8 is separate: it requires a linked bundle
at every required scale of every prescribed good point to obtain
the design class with its existence properties. Unfolding a given
$\mathrm{BWit}$ here uses the primitive rational tests and rows
already packed into its $\mathrm{Holds}$. To produce the original
selection one must construct those data upstream. This identifies
precisely where a new arithmetic or geometric existence proof enters.
% END B74R4


\paragraph{What remains for the base cusp space.}
On the original marked dictionary V74M3 supplies the vacuum
comparison used in J74T2. The remaining spanning input there is
stable generation by short sections on the identity cover. F74N3
supplies the level norm and K74A4 the reproducing integral. O74E5
gives the heat equations for the doubled theta--Eisenstein difference.
Its vanishing at zero and its annihilation by the invariant differential
operator, with the native transformation and convergence inputs,
are still needed for (ICL.6). To obtain the actual short-section
span one must also construct the common principal charts and the
remaining level data, and identify their analytic generators with the
algebraic sections as in (ICL.12). B74R4 applies once the common
linked selection exists.
\endgroup
